\documentclass[11pt,a4paper]{article}
\usepackage{fsproposal}
\usepackage[backend=biber,style=numeric-comp,sorting=none,maxbibnames=6,minbibnames=3,
            giveninits=true,doi=true,isbn=false,url=true,eprint=true]{biblatex}
\addbibresource{references.bib}
\AtEveryBibitem{\ifentrytype{article}{\clearfield{url}\clearfield{urldate}}{}}

\hypersetup{pdftitle={Experimental demonstration of reduced fibre solarisation in an oxygen-free atmosphere},
            pdfsubject={MSc thesis project proposal and student guide},
            pdfkeywords={fibre solarisation, UV, 369.5 nm, ytterbium ion, end-face degradation,
                         laser-induced contamination, oxygen absorber, 3D printing}}

\pagestyle{fancy}
\fancyhf{}
\fancyhead[L]{\sffamily\footnotesize\color{greydark}Reduced fibre solarisation in an oxygen-free atmosphere}
\fancyhead[R]{\sffamily\footnotesize\color{greydark}MSc proposal \& guide}
\fancyfoot[C]{\sffamily\footnotesize\thepage}
\renewcommand{\headrulewidth}{0.4pt}
\renewcommand{\headrule}{\hbox to\headwidth{\color{uvviolet!40}\leaders\hrule height \headrulewidth\hfill}}

\begin{document}

% =============================================================== title page
\begin{titlepage}
\thispagestyle{empty}
\begin{tikzpicture}[remember picture,overlay]
  \fill[uvviolet] (current page.north west) rectangle ([yshift=-5.2cm]current page.north east);
  \fill[uvviolet!70] ([yshift=-5.2cm]current page.north west) rectangle ([yshift=-5.45cm]current page.north east);
  \node[anchor=north west,text=white,font=\sffamily\large,inner sep=0] at ([xshift=2.4cm,yshift=-1.6cm]current page.north west)
    {MSc thesis project};
  \node[anchor=north west,text=white,font=\sffamily\Large\bfseries,inner sep=0] at ([xshift=2.4cm,yshift=-2.4cm]current page.north west)
    {Research proposal and student guide};
  \node[anchor=north west,text=white!85!uvviolet,font=\sffamily\small,inner sep=0] at ([xshift=2.4cm,yshift=-3.4cm]current page.north west)
    {Trapped-ion quantum technology \textbullet{} UV fibre delivery at 369.5\,nm};
\end{tikzpicture}
\vspace*{4.6cm}

{\raggedright\hyphenpenalty=10000\exhyphenpenalty=10000\sffamily\bfseries\fontsize{24}{29}\selectfont\color{uvviolet}
Experimental demonstration of reduced fibre solarisation in an oxygen-free atmosphere\par}
\vspace{6mm}
{\sffamily\large\color{greydark}
Can a 3D-printed controlled-atmosphere enclosure stop the UV-induced degradation of single-mode fibre
end faces carrying milliwatts of \SI{369.5}{nm} light?\par}
\vspace{12mm}

{\sffamily\small
\begin{tabular}{@{}>{\color{greydark}}l l@{}}
Student & \fillin{name} \\
Supervisor & \fillin{name} \\
Daily supervisor & \fillin{name} \\
Duration & 12 months full time (6-month variant in \cref{sec:workplan}) \\
Version & 1.0 -- 18 October 2026, draft for discussion \\
\end{tabular}\par}

\vfill
\begin{keybox}[Project at a glance]
\begin{itemize}[leftmargin=1.1em]
  \item \textbf{Problem.} Thorlabs polarisation-maintaining UV fibres carrying a few mW of
        \SI{369.5}{nm} light (the \Ybion{} cooling and detection wavelength) couple very well after a
        fresh polish, then lose throughput at a rate that grows with power, at both ends.
  \item \textbf{Idea.} Keep the end faces in an oxygen-free atmosphere, using 3D-printed cells with oxygen
        absorbers or an inert-gas purge, and show that the degradation stops.
  \item \textbf{Approach.} A multi-arm accelerated-ageing testbed runs ambient air, clean dry air, dry
        \Ntwo{} and sealed absorber cells \emph{in parallel} from one UV source. A fixed cleaning sequence
        and surface analysis locate the damage. The best solution is then demonstrated on a working
        \SI{369.5}{nm} link.
  \item \textbf{Why it is a good project.} Every outcome is useful. Either oxygen is the culprit and the
        fix is cheap, or the experiment identifies the real cause (contamination, water, particles or
        intrinsic glass damage) and what to do about it.
\end{itemize}
\end{keybox}
\end{titlepage}

\tableofcontents
\clearpage

\part*{Part I -- Research proposal}
\addcontentsline{toc}{part}{Part I -- Research proposal}
\section{Summary}\label{sec:summary}

Trapped-ion experiments with \Ybion{} deliver their \SI{369.5}{nm} cooling and detection light through
single-mode, polarisation-maintaining optical fibres. In the group, these fibres show a characteristic
failure. Directly after the connector end faces are polished, the coupling efficiency is excellent. It
then falls over days to weeks, faster at higher power, and both the input and the output end are
affected. Re-polishing restores the throughput, so the damage must sit at, or within micrometres of, the
end faces rather than in the bulk of the fibre. At a few milliwatts in a mode-field diameter of
about \SI{2.3}{\micro m}, the end faces see \SIrange{e5}{e6}{W/cm^2} of near-UV light, an
intensity at which photochemistry of the glass surface and of anything adsorbed on it becomes
important.

This project tests the hypothesis that \textbf{removing oxygen from the end faces halts the degradation},
and develops a practical way of doing so with \textbf{3D-printed enclosures} that are either purged with
inert gas or sealed with an \textbf{oxygen absorber}. The literature does not settle the question in advance. Oxygen
could drive the damage through reactions with photo-activated surface sites. In the laser-induced
contamination of UV optics in space lasers, however, oxygen is known to \emph{suppress} the growth of
absorbing deposits from organic vapours. Water vapour and airborne particles are further candidates.
The experiment is therefore designed to separate these factors. A multi-arm testbed exposes identical
fibres \emph{simultaneously} to ambient air, clean dry synthetic air, clean dry nitrogen and sealed
absorber cells, and later to controlled oxygen, humidity and contaminant levels. Throughput,
back-reflection and the atmosphere in every cell are logged continuously. At the end of each run a fixed
cleaning and re-polishing sequence, together with surface analysis, shows where the loss is and what
it consists of.

The project runs over 12 months (a 6-month variant is included). It starts with a forensic study of
fibres already degraded in the lab, builds and qualifies the testbed and cells, runs two atmosphere
campaigns, and ends with a demonstration on a working \SI{369.5}{nm} link. Success means a quantitative
answer to whether oxygen matters (with confidence intervals), and a qualified cell design that holds
$<\SI{100}{ppm}$ \Oxy{} for months. If the atmosphere is the key, the degradation should fall below
\SI{1}{\percent} per \SI{1000}{h} at \SI{3}{mW}. Part~II of this document is a practical guide for the
student: onboarding plan, core skills, working practices, pitfalls, a prioritised reading list and
draft procedures.

\section{Background and motivation}\label{sec:background}

\subsection{Near-UV light delivery in trapped-ion experiments}
In \Ybion{} experiments, light at \SI{369.5}{nm} drives the \ce{^2S_{1/2}}--\ce{^2P_{1/2}} transition used for
Doppler cooling, optical pumping and state detection. Light at \SI{399}{nm} is used for isotope-selective
photoionisation of neutral Yb. Delivering this light through single-mode fibres has clear advantages.
The fibre acts as a spatial filter that produces a clean Gaussian mode. It decouples the laser table
from the vacuum system and keeps beam pointing at the ion stable. With polarisation-maintaining (PM)
fibre, it also delivers a well-defined polarisation. Scalable architectures with many ion-trap modules
depend even more on fibre delivery, so the long-term reliability of near-UV fibres is a practical
concern for the whole field.

\subsection{The problem in the group}
The group uses Thorlabs PM patch cables designed for the near UV (referred to in the lab as
``375PM''; \cref{sec:sota-fibres}) at powers of a few milliwatts. The behaviour observed is:
\begin{itemize}
  \item directly after the end faces have been polished, the coupling efficiency is excellent;
  \item it then declines over time, faster when more power is used;
  \item both the input and the output end are affected;
  \item polishing the end faces again restores the throughput.
\end{itemize}
The practical consequences are frequent re-polishing (person-hours and risk of damaging the connector),
slow drifts of the power delivered to the ions, the need for spare laser power, and occasionally
sudden failures during measurements.

\paragraph{What these observations already tell us} Re-polishing removes only micrometres to tens of
micrometres of glass. Colour centres formed along the length of the fibre would survive it. The
restoration therefore places the damage at, or within micrometres of, the end faces. The power dependence
shows that the process is driven by light. The degradation of the \emph{output} face, which sees almost the
same guided intensity as the input face but no focusing optics, argues against simple misalignment of the
input coupling and for a process driven by local intensity at a glass--gas interface. At \SI{3}{mW} in a
\SI{2.3}{\micro m} mode-field diameter, the peak intensity is \SI{\approx1.4e5}{W/cm^2}, and each
square centimetre of the core region receives about \num{3e23} photons per second (\cref{app:estimates}).
At such photon fluxes even very weak absorption by surface sites or adsorbed molecules leads to
photochemistry within seconds.

\subsection{Terminology and scope}
In the lab, ``solarisation'' is used for any UV-induced loss of fibre throughput. Strictly, solarisation
means the formation of colour centres in the glass (\cref{sec:sota-bulk}). In this document,
\emph{end-face degradation} denotes the observed phenomenon and \emph{solarisation} denotes colour-centre
formation. Finding out which of the two is happening, or whether both are, is part of the project (RQ1).
The scope is the atmosphere at the end faces of the fibre type used in the group, at \SI{369.5}{nm} and
a nearby testbed wavelength (\SI{375}{nm}). Developing new fibres is out of scope, but alternative fibre
technologies are reviewed (\cref{sec:sota-fibres}) because they are the fall-back if the atmosphere turns
out not to be the cause.

\section{State of the art}\label{sec:sota}

\subsection{UV-induced defects in silica and their relevance at 369.5\,nm}\label{sec:sota-bulk}
\paragraph{Energetics} A \SI{369.5}{nm} photon carries \SI{3.36}{eV}; two photons carry \SI{6.71}{eV}.
The absorption edge of amorphous \ce{SiO2} lies near \SI{8}{eV}~\cite{Kajihara2010}, so neither one nor
two photons can excite the perfect glass across its gap. Excimer-laser studies show that E$'$ centres form
in a two-photon process at \SI{5.0}{eV} and \SI{6.4}{eV}~\cite{Arai1988,Tsai1988}. XeCl light
(\SI{4.0}{eV}, two-photon energy \SI{8.0}{eV}) created no centres in dry synthetic
silica~\cite{Arai1988}. At \SI{369.5}{nm} only processes that start at a pre-existing defect, an impurity
or a surface site are therefore expected: one-photon absorption into a weak band, or stepwise two-photon
absorption through a long-lived intermediate state.

\paragraph{The relevant defects} \Cref{tab:defects} lists the main point defects of silica
glass~\cite{Skuja1998,Griscom1991,Skuja2005,Skuja2012,Girard2019}. Most absorb well above
\SI{3.36}{eV}, but several channels exist at \SI{369.5}{nm}:
\begin{itemize}
  \item the spin-forbidden singlet--triplet band of the twofold-coordinated silicon centre ODC(II) at
        \SI{3.15}{eV}, whose long-lived (\si{ms}) triplet state can absorb a second photon;
  \item the \SI{6.9}{eV} band of ODC(II) and the onset of the Schumann--Runge absorption of interstitial
        \Oxy{} (shifted to $\approx\SI{6.5}{eV}$ in the glass), both reachable by two photons;
        photolysed interstitial \Oxy{} yields mobile oxygen atoms, interstitial ozone and peroxy
        radicals~\cite{Awazu1990,Kajihara2004a,Kajihara2005a,Skuja2000};
  \item chlorine species left by the manufacturing process: interstitial \ce{Cl2} absorbs at
        \SI{3.78}{eV} with a width of \SI{0.6}{eV}, and atomic \ce{Cl^0} at \SI{3.26}{eV}. In pure-silica-core
        fibres, much of the radiation-induced loss above \SI{2.5}{eV} is chlorine-related~\cite{Skuja2017,Girard2019}.
\end{itemize}
In pure-silica-core single-mode fibre under \SI{405}{nm} CW light at \si{MW/cm^2}, bleaching of ODC(II) and
growth of E$'$ centres have indeed been observed~\cite{Gonschior2013JNCS}. Germanium-related centres
would add a one-photon channel through the GeODC triplet band near \SI{3.8}{eV}. This is the origin of
the photosensitivity of Ge-doped fibres at near-UV and, via two photons, visible
wavelengths~\cite{Hill1978,Starodubov1997,Grubsky1999}. The group's PM-S350-HP fibre has a pure-silica core, so these
channels should be absent.

\begin{table}[t]
\centering\small
\caption{Main point defects in silica glass: absorption and photoluminescence (PL) bands, in eV, after
\cite{Skuja1998,Skuja2011,Skuja2017,Girard2019}. Values in bold are within reach of one (\SI{3.36}{eV})
or two (\SI{6.71}{eV}) photons at \SI{369.5}{nm}.}
\label{tab:defects}
\begin{tabularx}{\textwidth}{@{}L{4.6cm} L{3.6cm} L{2.4cm} Y@{}}
\toprule
Defect & Absorption (eV) & PL (eV) & Comment \\
\midrule
E$'$ centre, $\equiv$Si$^{\bullet}$ & 5.8 & -- & dominant UV-induced centre \\
NBOHC, $\equiv$Si--O$^{\bullet}$ & 1.97, 4.8 (continuum 4--7.8) & 1.91 & \SI{3.36}{eV} lies in its low-absorption gap \\
ODC(I), $\equiv$Si--Si$\equiv$ & 7.6 & 4.4 & precursor of E$'$ \\
ODC(II), $=$Si: & 5.05, \textbf{6.9}; triplet \textbf{3.15} & 4.4, 2.75 & triplet band and long-lived T$_1$ state \\
peroxy radical, $\equiv$Si--O--O$^{\bullet}$ & $\approx$2.0, 4.8 & -- & forms from NBOHC + O, or E$'$ + \Oxy{} \\
interstitial \Oxy{} & onset \textbf{$\approx$6.5} & 0.975 & source of O atoms and \ce{O3} \\
interstitial \ce{O3} & 4.8 & -- & bleached by UV \\
interstitial \ce{Cl2}, \ce{Cl^0} & \textbf{3.78} (FWHM 0.6), \textbf{3.26} & -- & chlorine from manufacture \\
GeODC (Ge-doped cores only) & 5.15; triplet \textbf{$\approx$3.8} & 4.3, 3.2 & Ge photosensitivity \\
\bottomrule
\end{tabularx}
\end{table}

\paragraph{Hydrogen, OH and fibre manufacture} Molecular hydrogen passivates E$'$ centres and NBOHCs. The
E$'$ + \Htwo{} reaction has an activation energy of only \SI{0.38}{eV} and proceeds at the same rate for
surface E$'$ centres~\cite{Radzig1995,Messina2007}. Under ArF and \ce{F2} irradiation \Htwo{} can also
promote ODC(I) formation~\cite{Ikuta2002}. Fibre drawing creates E$'$ centres, NBOHCs (the \SI{630}{nm}
band, which grows with draw tension) and peroxy radicals~\cite{Hanafusa1985,Hibino1986,Hibino1987}.
Strained Si--O--Si bonds are the main intrinsic precursors~\cite{Kajihara2008b}. Decades of work on
multimode UV fibres have led from hydrogen-loaded to hydrogen-free, high-OH designs with improved deep-UV
stability~\cite{Karlitschek1995,Karlitschek1998a,Klein2001,Khalilov2006,KleinKhalilov2015}. Loading with
\Htwo{} and UV curing has transferred this to single-mode PCF~\cite{Colombe2014}.

\paragraph{Kinetics} Colour-centre growth in fibres typically follows a stretched exponential or a
power law in dose~\cite{Griscom1993,Griscom2001}. Two-photon creation shows up as a quadratic
dependence on intensity~\cite{Arai1988,Tsai1988}. Comparable methods for measuring and fitting
photodarkening have been standardised for Yb-doped fibres~\cite{Koponen2006,Koponen2008,Jetschke2007}.
Measuring the initial rate at three or more powers therefore gives a cheap first indication of the
mechanism.

\paragraph{Gas transport into the glass} At room temperature \Htwo{} diffuses through a
\SI{125}{\micro m} fibre in a few days~\cite{Lemaire1991,Shang2009}. Oxygen molecules move about ten
orders of magnitude more slowly. Extrapolated high-temperature data~\cite{Norton1961,Kajihara2010} give
penetration depths of 0.1--5\,nm per year (\cref{app:estimates}). The bulk of the core is therefore
isolated from the atmosphere, except for deliberately added hydrogen. Oxygen can act only at the
surface, where it dissolves readily~\cite{Kajihara2004c}.

\paragraph{Reactive surfaces} Freshly fractured or ground silica carries reactive sites: silicon and
oxygen dangling bonds, silylene groups and strained two-membered rings~\cite{Radzig1993,Radzig1998,Radzig2000,Ceresoli2000}.
These react with \water{} (which opens the strained rings~\cite{Rimola2008}), with \Htwo{} and with
\Oxy{}. Polished end faces are such mechanically activated surfaces. Oxygen, water and hydrogen
therefore determine what UV-generated surface defects turn into: peroxy radicals, silanols or Si--H.
These products absorb mainly above \SI{4}{eV}, not at \SI{3.36}{eV}. Spectroscopically, H1 is therefore
not the most likely mechanism, which is one more reason to test it directly rather than assume it.


\subsection{What has been observed in single-mode fibres carrying UV and blue light}\label{sec:sota-reports}
\Cref{tab:reports} collects the most relevant reports. Two kinds of degradation emerge clearly.
\emph{Bulk solarisation} dominates in the deep UV. It is irreversible, it can set in at microwatt powers
in fused-silica fibres~\cite{Colombe2014}, and what light remains is carried in a speckled mode
pattern~\cite{Marciniak2017}. Below about \SI{300}{nm}, unprotected photonic-crystal fibres lose most of
their transmission within hours at milliwatt powers~\cite{Colombe2014,Gebert2014}. \emph{End-face
degradation} is reversible by re-cleaving or re-polishing, and it is the more relevant kind at near-UV and
blue wavelengths. It takes at least three forms:
\begin{itemize}
  \item \textbf{Deposits.} Colombe et al.\ mention UV-induced carbon deposition on input
        faces~\cite{Colombe2014}. Marciniak et al.\ saw deposits on large-mode-area PCF faces at
        \SI{313}{nm} with input powers as low as \SI{10}{mW}. Polishing removed them, and a slight
        overpressure of dry \Ntwo{} at the collimators prevented them for months~\cite{Marciniak2017}. They
        attributed the deposits to particles and oil drawn into the beam by the ``optical-tweezer effect''
        and then burned in. Purged and unpurged ends of the same fibre behaved differently. A Duke group
        reports the same use of nitrogen purging against particle deposition at UV
        wavelengths~\cite{YuJ2026}.
  \item \textbf{Restructuring of the silica.} In the most detailed study of CW blue light in small-core
        single-mode fibres (\SI{405}{nm}, mode-field diameter \SIrange{2}{3}{\micro m}), Gonschior et al.\
        found no deposits. Instead, laser-induced periodic surface structures formed on both faces:
        protrusions of \SIrange{0.5}{0.7}{\micro m} after 9--11 days, ripples of \SIrange{250}{400}{nm}
        period aligned with the polarisation, and a porous, lower-density silica layer. EDX and ToF-SIMS
        found no foreign elements~\cite{Gonschior2013JNCS,Gonschior2013PhD}. The loss grew faster than
        linearly in time, typically \SIrange{0.2}{1}{dB/day} at \SIrange{74}{95}{mW}. The output beam
        quality degraded ($M^2$ from 1.2 to 5 in 12 days). Cleaved faces degraded four times more slowly
        than polished ones~\cite{Gonschior2013PhD,Gonschior2014OE}. The structures appeared at
        \SIrange{0.5}{2.5}{MW/cm^2}, three orders of magnitude below the optical damage
        threshold~\cite{Gonschior2013JNCS}.
  \item \textbf{Slow drifts of unclear origin.} At \SI{369}{nm}, Sandia found lensed-fibre-to-chip coupling stable at
        \SIrange{7}{24}{mW} on short time scales but degrading over hours, with surface contamination of
        the fibre or chip facets among the suspected causes~\cite{Kwon2024}.
\end{itemize}
Two further observations bear directly on the oxygen hypothesis. Gonschior ran a PM single-mode fibre for
50 days at \SI{250}{mW} (\SI{405}{nm}) between \emph{nitrogen-purged, sealed} collimators. The fibre
lost only about \SI{1}{dB}, but cone-like structures still formed on both faces~\cite{Gonschior2013PhD}.
Inert gas alone therefore did not stop the restructuring at that (much higher) intensity. Second, in
ion-trap vacuum systems, \SIrange{369}{370}{nm} light caused losses in oxide mirror coatings that were
reversed by admitting oxygen or air~\cite{Gangloff2015,Ballance2017}. Absence of oxygen can therefore
\emph{harm} UV-illuminated oxide surfaces.

\begin{table}[t]
\centering\footnotesize
\caption{Selected reports on UV/blue degradation of single-mode fibres and related surfaces. PCF:
photonic-crystal fibre; LMA: large mode area; HC: hollow core.}
\label{tab:reports}
\begin{tabularx}{\textwidth}{@{}L{1.9cm} L{1.2cm} L{2.6cm} Y L{3.6cm}@{}}
\toprule
Ref. & $\lambda$ (nm) & Fibre / surface & Observation & Mitigation that worked \\
\midrule
\cite{Colombe2014} & 313, 280 & LMA PCF & bulk solarisation within hours at \si{mW}; carbon on input face
  & \Htwo{} loading + UV curing; collapsed, angled ends \\
\cite{Gebert2014} & 280 & solid-core PCF; kagom\'e HC & solid core $-$80\,\% in 3\,h at \SI{2}{mW}; HC unchanged after \SI{100}{h} at \SI{15}{mW}
  & hollow core \\
\cite{Yamamoto2009} & 250 & LMA PCF & damage mainly at the input; partly reversed by re-cleaving & low power ($<\SI{0.3}{mW}$) \\
\cite{Gonschior2013JNCS,Gonschior2013PhD} & 405 & SM and PM, MFD 2--3\,\textmu m & LIPSS and porous silica on both faces; $M^2$ grows; no foreign elements
  & end caps or launch fibre: $<\SI{0.03}{dB/day}$ \\
\cite{Marciniak2017} & 313 (369) & LMA-PM PCF & particle/oil deposits from \SI{10}{mW}; removed by polishing & dry \Ntwo{} overpressure \\
\cite{Kwon2024} & 369 & lensed fibre to chip & slow (hour-scale) coupling loss & -- \\
\cite{Gangloff2015,Ballance2017} & 369--422 & oxide mirror coatings in vacuum & UV-induced loss in vacuum & reversed by \Oxy{} or air \\
\cite{Hippler2016,Kokkinos2017} & 355 & optics with organic contaminants & deposit growth in vacuum & \Oxy{} partial pressure \\
\bottomrule
\end{tabularx}
\end{table}

\begin{keybox}[What is not known]
We have found no published measurement of the degradation of step-index, pure-silica-core PM fibre at
\SI{369.5}{nm} and a few milliwatts, the configuration used in the group. No study in the trapped-ion
community has treated \Oxy{} as a controlled variable. Nitrogen purging is used to keep particles away,
not to remove oxygen. This is the gap the project fills.
\end{keybox}

\subsection{Fibre technologies for near-UV delivery}\label{sec:sota-fibres}
\paragraph{The group's fibre} The cables known in the group as ``375PM'' are Thorlabs P3-375PM-FC
patch cables (\SIrange{1}{10}{m}), FC/APC at both ends, aligned at \SI{375}{nm}, with no AR-coated
option~\cite{ThorlabsP3375PM}. They are built on PM-S350-HP: a PANDA fibre with a \emph{pure-silica
core}, mode-field diameter \SI{2.3}{\micro m} at \SI{350}{nm} (about \SI{2.4}{\micro m} at
\SI{369.5}{nm}), NA~0.12 and second-mode cut-off $\leq\SI{340}{nm}$~\cite{ThorlabsPMS350}. A
germanium-free core avoids the strong photosensitivity of Ge-doped silica, and Thorlabs advertises
pure-silica cores for their resistance to photodarkening. The same catalogue nevertheless states that all
fibres eventually photodarken in the UV and should be treated as consumables. For CW light at an
air--glass face of undoped silica, Thorlabs quotes a practical safe intensity of about
\SI{250}{kW/cm^2}~\cite{ThorlabsDamage}. With Thorlabs' convention this corresponds to about \SI{11}{mW}
for a \SI{2.4}{\micro m} mode. The group's few milliwatts are below that, yet the faces still degrade,
which illustrates that ``damage threshold'' figures say little about slow photochemical processes.
Thorlabs no longer sells the bare fibre; Coherent (Nufern) offers a fibre with the same part number. The
\SI{405}{nm} PM cables use PM-S405-XP, whose cut-off ($380\pm\SI{20}{nm}$) may allow a second mode at
\SI{369.5}{nm}. They are therefore no substitute.

\paragraph{Large-mode-area photonic-crystal fibre with hydrogen loading} Endlessly single-mode PCFs
spread the light over roughly ten times the mode area. Loading them with \Htwo{} (about \SI{10}{MPa},
room temperature, 4--6 days) and curing them with UV light before the hydrogen escapes makes them resistant
to deep-UV solarisation for hundreds of hours at \SI{125}{mW} (\SI{313}{nm})~\cite{Colombe2014,NISTrecipe}.
Hydrogen loading alone does not confer resistance unless curing occurs before the hydrogen has diffused
out~\cite{Colombe2014}. Collapsing the holes over a few hundred micrometres and angle-cleaving the ends
protects the faces and suppresses back-reflections. PM versions reach polarisation extinction ratios of
30--32\,dB at \SI{369}{nm} after curing~\cite{Marciniak2017}. These fibres are not off-the-shelf
patch cords, and their faces still collect deposits without a purge~\cite{Marciniak2017}.

\paragraph{Hollow-core fibres} Kagom\'e and antiresonant hollow-core fibres guide almost all the light in
gas, so the glass hardly sees it. They have transmitted \SI{280}{nm} CW light without change for
\SI{100}{h}~\cite{Gebert2014} and \SI{320}{nm} light at up to \SI{60}{mW} without degradation, at about
\SI{0.13}{dB/m}~\cite{Dorer2024}. Losses below \SI{10}{dB/km} have been reported at
\SI{369}{nm}~\cite{Osorio2023}, and guidance extends down to \SI{190}{nm}~\cite{Mears2024}. Caveats:
there is no PM hollow-core fibre, ammonium chloride can grow on cleaved hollow-core end
faces~\cite{Rikimi2020}, and reversible photodarkening has been reported in UV-guiding hollow-core
fibres~\cite{Mears2024SPIE}.

\paragraph{End caps} A coreless silica cap spliced to the fibre lets the beam expand before it reaches the
glass--gas interface. At \SI{369.5}{nm} a \SI{300}{\micro m} cap reduces the face intensity about
280-fold (\cref{app:estimates}). At \SI{405}{nm}, a short launch fibre and end caps reduced the
damage rate from \SIrange{0.2}{1}{dB/day} to below \SI{0.03}{dB/day}~\cite{Gonschior2013PhD}. End caps
attack the intensity-driven surface processes at their root. Together with the atmosphere control
studied here they are the most promising engineering solution, and the obvious next step if the
atmosphere alone is not enough.

\subsection{Surface processes and the role of the atmosphere}\label{sec:sota-surface}
Four communities have studied UV-driven degradation of optical surfaces in controlled atmospheres:
space lidar, UV frequency-conversion cavities, semiconductor lasers and laser modules, and
high-power fibre connectors. Their results are the best available guide to what the atmosphere around a
fibre end face will do.

\paragraph{Laser-induced contamination in space lasers} A prototype of the \SI{355}{nm} ALADIN laser for
ESA's Aeolus mission lost \SI{50}{\percent} of its transmission in less than six hours. Outgassed organics
were carbonised on the optics by the UV light, and even a few parts per billion of organics caused
contamination~\cite{ESA2018Aeolus}. The flight laser therefore operates with a low-pressure oxygen purge
of a few tens of pascals~\cite{Wernham2010,ESA2018Aeolus}. Systematic tests at DLR showed that an
\Oxy-to-contaminant pressure ratio above 50:1 drastically reduces deposit growth, above 1000:1 no
transmission loss was detectable, and existing deposits were completely removed by irradiation in
several tens of pascals of oxygen~\cite{Hippler2016,Tighe2008,Riede2011,Schroeder2013}. The ESA/DLR
guidance for ground operation is explicit: ``operation with normal dry air is to be preferred to operation
in nitrogen''~\cite{Mondin2017}. The same paper notes that silicones are not removed by oxidation, that CW
irradiation can produce more deposit than pulsed irradiation, and that high-purity bottled gas alone does
not prevent absorbing ``black dots''. In sealed \SI{1064}{nm} lasers, trace contamination from
ultra-high-purity nitrogen cylinders was enough to cause damage above about
\SI{60}{\celsius}~\cite{Hovis1994}. With \SI{300}{ppm} toluene
present, optics survived $8\times10^3$ pulses in \Ntwo, $1.8\times10^4$ in synthetic air, and more than
$10^6$ in air with a few thousand ppm of water or alcohol~\cite{Weiller2010}.

\paragraph{CW evidence at our intensity} Parthasarathy et al.\ studied contamination at a diamond
surface in vacuum under \SI{1.5}{mW} of CW \SI{532}{nm} light focused to about \SI{220}{kW/cm^2}, an
intensity close to that on the group's fibre faces~\cite{Parthasarathy2024}. Deposits grew
\emph{fastest} at an oxygen partial pressure of about \SI{5e-3}{mbar}. Between \SI{2e-3}{mbar} and
\SI{2}{mbar} their onset was delayed, and above \SI{2}{mbar} none formed. Humidity slowed the growth. At
atmospheric pressure, \SIrange{1}{1000}{ppm} of \Oxy{} corresponds to \SIrange{0.001}{1}{mbar}. An
``almost oxygen-free'' cell is therefore exactly in the regime where contamination was promoted, whereas
air is a hundred times above the inhibition threshold. How these vacuum results transfer to
\SI{1}{atm} of nitrogen is not known. That is one of the questions the \arm{O2-x} arms answer.

\paragraph{UV conversion cavities} Groups producing CW UV light for trapped ions routinely flush their
nonlinear crystals with oxygen. At NIST the \SI{313}{nm} BBO crystal sits in an enclosure ``purged with a
gentle flow of dry oxygen'', there mainly against moisture-induced fogging~\cite{Wilson2011}. At PTB oxygen is led directly onto the crystal facets ``to
prevent light-induced damage''~\cite{Hannig2018}. At \SI{243}{nm}, cavity mirrors degraded within minutes
below about \SI{0.7}{mbar} of oxygen, and the deposit could be removed with organic
solvents~\cite{Cooper2018,Burkley2021}. A recent watt-level \SI{230}{nm} source uses filtered dry air
because ``pure nitrogen removes this cleaning pathway''~\cite{Cai2026}. One industrial patent reports the
opposite for BBO: clean dry air alone degraded the crystal within 100 hours, and adding about \SI{1000}{ppm}
of water to the purge cured this~\cite{KLApatent}.

\paragraph{Semiconductor lasers and fibre-coupled laser modules} GaN laser diodes operated in pure
nitrogen grow a thick, highly absorbing deposit on the facet, which oxygen removes~\cite{Schoedl2005}.
Carbon deposition~\cite{Kim2008} and the growth of \ce{SiO_x} layers~\cite{Freier2024,WangXW2021} are
documented facet failure modes. Package patents mostly specify oxygen-containing fills: more than
\SI{100}{ppm}~\cite{Corningpatent}, at least \SI{5}{\percent}~\cite{Sanyopatent}, or
\SIrange{1}{18}{\percent}~\cite{Sharppatent}. For \SI{405}{nm} lasers, oxygen is said to react with
photo-decomposed silicon compounds before they deposit on the facet~\cite{Panasonicpatent}. The important
exception concerns \emph{fibre} end faces. In a fibre-coupled \SI{410}{nm} module, reliability
\emph{deteriorated} at oxygen concentrations of \SI{100}{ppm} or more, because silicon-based compounds
deposited at the fibre end face. Silicone parts were the principal source, and the patent specifies
\SIrange{1}{100}{ppm} \Oxy~\cite{Fujifilmpatent2006}. A related patent on siloxane-derived \ce{SiO_x}
deposits at a fibre input recommends silicone-free materials and air admitted only through activated
carbon~\cite{Ricohpatent}. Another describes a ``dust collection effect'' where the optical density is
high~\cite{Fujifilmpatent2007}.

\paragraph{Can 369.5\,nm light make ozone?} Splitting \Oxy{} needs photons below \SI{242}{nm} (the
Herzberg continuum~\cite{Yoshino1988}). UV/ozone cleaning therefore uses the \SI{185}{nm} and
\SI{254}{nm} mercury lines~\cite{Vig1985}. Two \SI{369.5}{nm} photons carry enough energy, but no
two-photon cross-section near \SI{370}{nm} is known, and any effect would be confined to the
\SI{\approx20}{\micro m} next to the face. Oxygen suppresses contamination even where it cannot be
photodissociated~\cite{Parthasarathy2024,Hippler2016}. The relevant chemistry is light-assisted oxidation
at the surface, not gas-phase ozone.

\paragraph{Particles and high-power connectors} Dielectric particles are trapped by the optical gradient
force in air~\cite{Ashkin1986,Omori1997}, and absorbing particles are moved by photophoretic
forces~\cite{Shvedov2009}. Fibre connectors carrying watts at \SI{1550}{nm} fail mainly because of
absorbing contamination on the face, at intensities comparable to ours~\cite{DeRosa2002,Seo2003,Domingues2011}.
The \SI{313}{nm} fibre deposits of Marciniak et al.\ were attributed to this mechanism~\cite{Marciniak2017}.

\begin{keybox}[Oxygen: friend or foe?]
\begin{itemize}[leftmargin=1em]
  \item \textbf{Hydrocarbon contamination:} every community finds oxygen (and often water vapour)
        \emph{protective}. Nitrogen and vacuum favour absorbing carbon deposits, most strongly at
        trace \Oxy{} levels.
  \item \textbf{Siloxane contamination:} oxidation turns silicones into non-volatile \ce{SiO_x}. At a
        \SI{410}{nm} fibre end face, oxygen above \SI{100}{ppm} \emph{accelerated} such deposits. This is
        the strongest support in the literature for the group's hypothesis.
  \item \textbf{Particles and silica restructuring:} not expected to depend on oxygen. A clean, slightly
        pressurised gas keeps particles away.
\end{itemize}
The outcome for the group's fibres therefore depends on which contaminants are present, and finding
that out is the first task (Phase~0). A sealed, nearly oxygen-free cell that contains outgassing
polymers is a high-risk configuration. The experiment tests it against clean, oxygen-containing gas
rather than assuming it helps.
\end{keybox}


\subsection{Creating and verifying a clean, oxygen-free atmosphere}\label{sec:sota-atmosphere}
\paragraph{Oxygen absorbers} Iron-based absorbers (Mitsubishi Gas Chemical AGELESS) reduce the oxygen
in a sealed barrier container to $\leq\SI{0.1}{\percent}$. The container wall should let through no more
than \SI{20}{mL} of oxygen per square metre, day and atmosphere of pressure difference~\cite{MGCAgeless}. The number in the type name is the volume of
\Oxy{} absorbed: a ZP-50 binds \SI{50}{mL}. Some types need moisture from their surroundings. The
self-reacting Z types work best above about \SI{55}{\percent} RH and release humidity into the
container~\cite{Daniel1993}. The reaction is exothermic, and the sachets lose capacity within minutes in
air. Lambert et al.\ give a useful rule for a leaking box: the equilibrium concentration is
$[\Oxy]\approx L/(12.7\,C)$, where $L$ is the ingress in ppm per day and $C$ is the absorber capacity
divided by the box volume~\cite{Lambert1992}. The RP System of the same manufacturer exists in two
types. RP-A absorbs oxygen, water and corrosive gases and drives the humidity to nearly zero. RP-K
leaves the humidity unchanged but can release organic acids if overloaded~\cite{MGCRP}. Both reach
$\leq\SI{0.1}{\percent}$ \Oxy{} within a day and shed fine powder, so the sachet belongs behind a
membrane. Colour indicators change near \SIrange{0.1}{0.5}{\percent} \Oxy{} and show only gross
failures. Vacuum non-evaporable getters need activation at \SIrange{350}{450}{\celsius} and would
saturate almost immediately in a cell at atmospheric pressure~\cite{SAESSorbAC}. In heated form they serve
as in-line gas purifiers. Commercial point-of-use purifiers deliver inert gas with
$<\SI{100}{ppt}$ \Oxy{} and \water{}~\cite{EntegrisHX}. A palladium catalyst with a few per cent \Htwo{}
holds anaerobic chambers at \SIrange{0}{5}{ppm} \Oxy~\cite{Coy}, an option if forming gas is used anyway.

\paragraph{Measuring oxygen, water and organic vapour} Proving $<\SI{100}{ppm}$ needs an instrument that
resolves ppm.
\begin{itemize}
  \item Ultra-trace luminescent sensor spots read through a window: e.g.\ PreSens SP-PSt9,
        \SIrange{0}{200}{ppm}, detection limit \SI{0.5}{ppm}~\cite{PreSensPSt9}. They are ideal for
        sealed cells. The common trace spots with a \SI{50}{ppm} detection limit are only marginal for this
        purpose~\cite{PyroTROX}, and percent-level sensors are useless here~\cite{LuminOx}.
  \item Galvanic trace analysers (detection limit \SI{\sim0.05}{ppm}) need a sample flow of
        \SIrange{0.5}{1}{L/min}. They suit the exhaust of purged arms rather than sealed
        cells~\cite{GPR1200}.
  \item Zirconia sensors read wrongly in hydrogen-containing gas and are poisoned by silicone
        vapour~\cite{PSTzirconia}.
\end{itemize}
Luminescent spots must be shielded from the UV beam. Capacitive humidity sensors resolve trends down to
about \SI{1}{\percent} RH (\SI{\approx230}{ppmv})~\cite{SHT4x}. Low-cost metal-oxide VOC sensors depend on
oxygen adsorbed on the oxide~\cite{Barsan2001}, so they work only in oxygen-containing streams. Organic
vapour in oxygen-free streams is best sampled onto sorbent tubes and analysed by TD-GC/MS.

\paragraph{Gas purity} Nitrogen 5.0 contains up to \SI{3}{ppm} \Oxy, \SI{5}{ppm} \water{} and
\SI{0.2}{ppm} hydrocarbons; 6.0 grade one sixth of that~\cite{LindeN2}. ``Zero grade'' air may contain
\SI{5}{ppm} of hydrocarbons~\cite{BOCzero}, whereas ultra-zero grades are specified at
\SI{0.1}{ppm}~\cite{AirgasUZ}. The hydrocarbon content must be specified explicitly, because
parts-per-billion levels already matter (\cref{app:estimates}).

\paragraph{3D-printed enclosures} FDM parts become leak-tight with slight over-extrusion (extrusion
multiplier 0.98 to 1.2 depending on the printer), walls of at least \SI{1.1}{mm} and randomised
seams~\cite{Gordeev2018}. Thin layers matter more than infill~\cite{Leeder2025}, and laser-sintered PA12
permeates like bulk material~\cite{Liebrich2021}. Printed parts are used in vacuum systems: SLS PA12 was
characterised by Povilus et al.~\cite{Povilus2014}; FDM PEEK, SLA resins and other printed plastics have
been tested since~\cite{Stallwood2018,Radic2025,Domingues2025}. Printed metal chambers reach ultra-high
vacuum in portable quantum technology~\cite{Cooper2021}. \Cref{tab:materials} lists permeabilities and
outgassing data. Three points stand out:
\begin{itemize}
  \item Dense polymer walls let in 50--500 times more water molecules than oxygen molecules. Printed ABS
        stores \SI{0.35}{wt\percent} of gas, mostly water~\cite{Sefa2016}. A sealed ``dry'' cell
        therefore becomes wet within hours unless it contains a desiccant (\cref{app:estimates}).
  \item Standard SLA resins lose \SIrange{1.5}{2.9}{\percent} of their mass in vacuum, whereas PETG,
        PEEK, ULTEM and Antero lose about \SI{0.3}{\percent}~\cite{NASAoutgassing}. FDM printing of
        ABS releases styrene, and cured resins still emit acrylates and
        aldehydes~\cite{Azimi2016,Wojnowski2022,Pham2023}.
  \item O-rings need a surface finish of about \SI{0.4}{\micro m}, squeeze of at most \SI{30}{\percent}
        and \SIrange{60}{85}{\percent} gland fill~\cite{ParkerORing}. FDM surfaces are much rougher,
        so the seal should sit on an ironed top skin, a machined land or an SLA part. Silicone O-rings
        let in about 200 times more oxygen than FKM or butyl, and they release siloxanes.
\end{itemize}

\begin{table}[t]
\centering\small
\caption{Permeability at about \SI{23}{\celsius} (1\,Barrer $=\SI{3.35e-16}{mol.m.m^{-2}.s^{-1}.Pa^{-1}}$)
and outgassing (ASTM E595 total mass loss TML and collected volatile condensable material CVCM) of
candidate cell materials. Values compiled from \cite{Marano2022,Levitsky2019,Sonstevold2023,Liebrich2021,Peacock1980,Chiggiato2020,NASAoutgassing};
-- not found. Permeabilities of PETG, ASA, ABS and PEI should be measured by rate of rise.}
\label{tab:materials}
\footnotesize
\begin{tabularx}{\textwidth}{@{}l c c c Y@{}}
\toprule
Material & \Oxy{} (Barrer) & \water{} (Barrer) & TML / CVCM (\%) & Remark \\
\midrule
PET / PETG          & 0.034 (PET) & 130 (PET) & 0.35 / 0.02 (PETG) & good barrier; low VOC emission when printed \\
PLA                 & 0.15--0.43  & 2000--3300 & 0.70--0.75 / 0.01 & very water-permeable \\
ABS                 & --          & $\approx$2600 & 0.59--0.94 / 0.01--0.04 & emits styrene; avoid \\
PC                  & 1.0         & --         & 0.14--0.17 / 0.00 & \\
PP                  & 1.7         & --         & -- & poor \Oxy{} barrier \\
PA12 (SLS)          & 0.13--0.20  & --         & 0.75 / 0.02 & hygroscopic \\
PEEK                & 0.14        & $\approx$140 & 0.26--0.30 / $\leq$0.02 & needs a high-temp.\ printer \\
ULTEM 9085 / Antero 800NA & --    & --         & 0.33--0.37 / $\leq$0.02; 0.35 / 0.004 & low-outgassing FDM \\
SLA resins (Clear, Tough, Grey) & 0.125 (VeroClear) & -- & 1.5--2.9 / $\leq$0.02 & Rigid 10K: 0.54 / 0.00 \\
PVDF / PCTFE        & 0.031 / 0.093 & --       & 0.27 / 0.01 (PVDF) & excellent barriers \\
PTFE / FEP          & 2.3 / 4.5   & --         & -- & poor \Oxy{} barrier despite inertness \\
FKM / butyl O-rings & 1.3--1.7    & 40--200    & 0.10--0.82 / $\leq$0.04 (FKM) & preferred seal materials \\
polyurethane        & 1.5--4.7    & 340--12\,500 & -- & avoid for gas lines \\
silicone            & 100--600    & 10\,500    & -- & never near the end faces \\
\bottomrule
\end{tabularx}
\end{table}


\subsection{Technical facts that shape the testbed}\label{sec:sota-hardware}
\begin{itemize}
  \item \textbf{Source.} The only single-mode \SI{375}{nm} laser diode in Nichia's line-up is the
        NDU4116 (\SI{70}{mW}), sold by Thorlabs as L375P70MLD~\cite{ThorlabsL375}. It needs a TEC
        mount that supports its pin code (e.g.\ LDM56F). Splitting its free-space beam before any fibre
        makes every input face a test item, and no upstream delivery fibre ages. With \SIrange{30}{60}{\percent}
        coupling efficiency, four to six arms at \SI{3}{mW} guided are possible.
  \item \textbf{Detectors.} Silicon photodiodes lose UV responsivity under exposure~\cite{Korde1987,Siles2024};
        photodiodes with high UV resistance exist~\cite{HamamatsuS12698}. SiC photodiodes respond only weakly
        at \SI{369.5}{nm} because it lies at their band edge.
  \item \textbf{Lenses.} The moulded aspheres commonly used for fibre coupling are made of high-index
        glasses whose UV transmission is reduced. Check them at \SI{369.5}{nm}.
  \item \textbf{Inspection standards.} The core zone of IEC~61300-3-35 has a diameter of
        \SI{25}{\micro m}, about ten times the mode-field diameter~\cite{IEC61300335}. Standard
        pass/fail criteria are too coarse for these fibres. An in-house criterion based on
        $\geq400\times$ imaging of the central \SI{10}{\micro m} is needed.
  \item \textbf{Polishing.} Thorlabs' procedure for FC/APC uses diamond films down to \SI{1}{\micro m}
        and a final \SI{0.01}{\micro m} silica film, and warns that alumina and SiC films undercut fibres
        in zirconia ferrules~\cite{ThorlabsPolish}. Every final film and lubricant also leaves a
        possible residue on the face. The procedure is therefore part of the experimental protocol.
\end{itemize}

\subsection{Contribution of this project}
The project combines four elements that, to our knowledge, have not been brought together:
(i) the step-index pure-silica PM fibre and the wavelength and power actually used in \Ybion{}
experiments; (ii) atmosphere as a controlled variable, with \Oxy, water, organic vapour and particles
separated; (iii) simultaneous, identical arms with continuous logging of throughput, face reflection and
gas composition; and (iv) a low-cost, deployable enclosure qualified against the same measurements.

\section{Aim, research questions and hypotheses}\label{sec:aims}

\subsection{Aim}
The aim of the project is to \textbf{determine whether, and under which atmospheric conditions, the
UV-induced degradation of single-mode fibre end faces carrying a few milliwatts of \SI{369.5}{nm}
light can be stopped, and to demonstrate a low-cost, 3D-printed end-face enclosure that achieves this
in routine use}. The working hypothesis of the group is that removing oxygen from the end faces halts
the degradation. The project is designed so that it gives a clear answer whether this
hypothesis is right or wrong. If oxygen turns out not to be the culprit, the same
measurements identify what is (\cref{tab:predictions}).

\subsection{Research questions}
\begin{enumerate}[label=\textbf{RQ\arabic*},leftmargin=3.2em]
  \item \textbf{Where and what?} Is the throughput loss located at the input face, the output face or in
        the fibre, and what is its physical nature (deposit, surface damage, colour centres, or
        misalignment)?
  \item \textbf{Is oxygen the cause?} At fixed (low) humidity and contaminant levels, does removing
        \Oxy{} change the degradation rate, by how much, and is there a threshold \Oxy{}
        concentration below which the rate is negligible?
  \item \textbf{What else matters?} What are the separate roles of water vapour, organic and siloxane
        vapours (including those emitted by 3D-printed parts, adhesives and cable jackets), airborne
        particles and, optionally, hydrogen?
  \item \textbf{Does it work in practice?} Can a passive, sealed 3D-printed cell containing an oxygen
        absorber hold the required atmosphere for months, and does it stop the degradation of a
        working \SI{369.5}{nm} fibre link in one of the group's experiments?
\end{enumerate}

\subsection{Competing hypotheses}
The observations available so far (high throughput after re-polishing, decline at a rate that
increases with power, decline at \emph{both} ends) constrain the possible mechanisms considerably. Because
re-polishing removes only a few micrometres to tens of micrometres of glass and fully restores the
throughput, the damage must sit at, or within micrometres of, the end faces. Bulk colour-centre
formation along the fibre would not be undone by polishing. The guided intensity is nearly the same at
both faces, so a process driven by the local optical intensity at a glass surface explains why the output
face degrades too. Mechanical misalignment of the input coupling does not. Six hypotheses are retained
(\cref{fig:mechanisms}):

\begin{description}[leftmargin=2.6em,style=nextline,itemsep=3pt]
  \item[H1 -- oxygen-driven surface photochemistry] UV light (most likely via two-photon excitation at
        \SI{6.7}{eV}) creates reactive sites on the polished silica surface. Gas-phase \Oxy{} reacts with them
        to form absorbing or scattering oxygen-excess defects or oxidised layers. \emph{Prediction:}
        removing \Oxy{} stops the degradation, whatever the humidity or contaminant level.
  \item[H2 -- laser-induced contamination (LIC)] Molecules adsorbed from the gas phase (hydrocarbons,
        siloxanes, plasticisers) are photo-polymerised or photo-decomposed into an absorbing film that
        grows with dose. In the LIC literature oxygen \emph{suppresses} deposit growth by photo-oxidising
        the deposit~\cite{Riede2011,Kokkinos2017}. \emph{Prediction:} clean gas slows degradation. Removing \Oxy{}
        while organics remain (e.g.\ outgassing from printed parts) makes it \emph{worse}.
  \item[H3 -- water-assisted surface damage] Adsorbed water takes part in the surface photochemistry
        (hydroxylation, etching, roughening). \emph{Prediction:} dry gas helps whether or not
        \Oxy{} is present. Humidified \Ntwo{} degrades the face.
  \item[H4 -- particle accumulation] Airborne particles are pulled into the high-intensity spot by
        optical and photophoretic forces and then burn in. \emph{Prediction:} any filtered gas helps,
        and particles are visible on the end face.
  \item[H5 -- intrinsic glass damage] Colour centres form in the core within the polished-away depth
        (for example a defect-rich layer left by polishing). \emph{Prediction:} no atmosphere effect,
        except a possible reduction in hydrogen-containing gas~\cite{Colombe2014}.
  \item[H6 -- mechanical or thermal artefact] Drift of the coupling optics, epoxy creep or ferrule movement.
        \emph{Prediction:} no atmosphere effect, and throughput is recovered by re-peaking the alignment without
        touching the fibre.
\end{description}
Several mechanisms may act together. In particular, H1 and H2 are not exclusive. Any absorbing film
also heats the end face (\cref{app:estimates}), which can turn a slow process into a runaway one. The
experiment is therefore designed to measure the effect of each atmospheric factor separately, not just
to compare ``with'' and ``without'' oxygen.

\begin{table}[t]
\centering\small
\caption{Predicted degradation rate in each test atmosphere relative to the ambient-air control, for
each hypothesis. \arm{AIR}: ambient lab air; \arm{CDA}: clean, dry synthetic air; \arm{N2}: clean, dry
nitrogen; \arm{ABS}: sealed printed cell with \Oxy{} absorber; \arm{WET-N2}: humidified \Ntwo{};
\arm{FG}: forming gas (\SI{5}{\percent} \Htwo{} in \Ntwo{}); \arm{N2+P}: nitrogen with a sample of printing
polymer in the cell. Symbols: $\bullet$ similar to \arm{AIR}; $\circ$ strongly reduced;
$\uparrow$ faster than \arm{AIR}; -- no clear prediction. The comparison \arm{CDA}~vs~\arm{N2} is the
decisive test of H1 versus H2.}
\label{tab:predictions}
\setlength{\tabcolsep}{5pt}
\begin{tabular}{@{}l ccccccc l@{}}
\toprule
 & \arm{AIR} & \arm{CDA} & \arm{N2} & \arm{ABS} & \arm{WET-N2} & \arm{FG} & \arm{N2+P} & extra signature \\
\midrule
H1 \Oxy{} photochemistry   & $\bullet$ & $\bullet$ & $\circ$   & $\circ$         & $\circ$   & $\circ$   & $\circ$     & rate rises with \Oxy{} fraction \\
H2 contamination (LIC)     & $\bullet$ & $\circ$   & $\bullet$/$\circ$ & $\bullet$/$\uparrow$ & --  & --        & $\uparrow$  & C/Si-rich film; removed by \Oxy{} plasma \\
H3 water                   & $\bullet$ & $\circ$   & $\circ$   & $\circ$         & $\bullet$ & $\circ$   & $\circ$     & correlates with lab RH \\
H4 particles               & $\bullet$ & $\circ$   & $\circ$   & $\circ$         & $\circ$   & $\circ$   & $\circ$     & discrete particles on core \\
H5 glass damage            & $\bullet$ & $\bullet$ & $\bullet$ & $\bullet$       & $\bullet$ & $\circ$?  & $\bullet$   & no surface film; PL bands \\
H6 mechanical/thermal      & $\bullet$ & $\bullet$ & $\bullet$ & $\bullet$       & $\bullet$ & $\bullet$ & $\bullet$   & recovered by re-peaking \\
\bottomrule
\end{tabular}
\end{table}

\begin{keybox}[Why oxygen can only act at the surface]
At room temperature molecular oxygen is practically immobile in silica glass. With literature diffusion
data it penetrates far less than a nanometre in a year (\cref{app:estimates}), so the external
atmosphere cannot change bulk solarisation of the core. Hydrogen is different: it diffuses through a
\SI{125}{\micro m} fibre within days~\cite{Lemaire1991}. Any oxygen effect must therefore act at the end
faces. The observed behaviour of the group's fibres (restored by re-polishing, both ends affected) points
to exactly such surface processes, which makes the project well motivated.
\end{keybox}


\begin{figure}[p]
  \centering
  % Candidate degradation processes at a UV-illuminated fibre end face (input end shown).
\begin{tikzpicture}[font=\sffamily\scriptsize,>=Stealth,x=1cm,y=1cm]
  % --- ferrule, cladding, core
  \fill[gray!22] (0,-1.25) rectangle (6.6,1.25);
  \fill[cyan!7] (0,-0.55) rectangle (6.6,0.55);
  \draw[cyan!45!gray,thin] (0,-0.55) -- (6.6,-0.55) (0,0.55) -- (6.6,0.55);
  \fill[uvviolet!40] (0,-0.11) rectangle (6.6,0.11);
  \draw[gray!65!black,thick] (6.6,-1.25) -- (6.6,1.25);
  \draw[gray!55] (0,-1.25) -- (6.6,-1.25) (0,1.25) -- (6.6,1.25);
  \node[gray!50!black,anchor=west] at (0.1,0.92) {zirconia ferrule};
  \node[gray!50!black,anchor=west] at (0.1,0.36) {silica cladding};
  \node[white,anchor=west] at (0.1,0) {\textbf{core}};
  \draw[decorate,decoration={snake,amplitude=0.6mm,segment length=3mm},gray!60] (-0.02,-1.25) -- (-0.02,1.25);
  % --- converging UV beam
  \fill[uvviolet!22] (6.6,0.10) -- (10.9,1.2) -- (10.9,-1.2) -- (6.6,-0.10) -- cycle;
  \draw[uvviolet!70,thin] (6.6,0.10) -- (10.9,1.2) (6.6,-0.10) -- (10.9,-1.2);
  \draw[->,uvviolet,very thick] (10.75,0) -- (8.9,0);
  \node[uvviolet!80!black,anchor=west] at (9.0,0.24) {\textbf{369.5\,nm}};
  \node[uvviolet!80!black,anchor=west] at (9.0,-0.26) {$\sim\!10^{5}$\,W\,cm$^{-2}$};
  % --- H2: deposit on the core region
  \fill[amberdark!70] (6.6,-0.30) .. controls (6.74,-0.22) and (6.76,0.22) .. (6.6,0.30) -- cycle;
  % --- H5: colour centres in the glass
  \foreach \x in {4.7,5.25,5.7,6.05,6.32}{\node[reddark,inner sep=0] at (\x,0.0) {$\ast$};}
  % --- H1: O2 molecules
  \foreach \x/\y in {7.35/1.55,8.05/1.85,7.0/1.95}{
    \begin{scope}[shift={(\x,\y)}]
      \fill[tealdark] (0,0) circle (0.07); \fill[tealdark] (0.13,0) circle (0.07);
    \end{scope}}
  \draw[->,tealdark,thin] (7.3,1.45) .. controls (6.95,1.1) .. (6.67,0.2);
  % --- H3: adsorbed water
  \foreach \x/\y in {7.25/-1.65,7.9/-1.9}{
    \begin{scope}[shift={(\x,\y)}]
      \fill[cyan!60!blue] (0,0) circle (0.075);
      \fill[gray!60] (-0.09,0.07) circle (0.04); \fill[gray!60] (0.09,0.07) circle (0.04);
    \end{scope}}
  \draw[->,cyan!60!blue,thin] (7.2,-1.55) .. controls (6.9,-1.1) .. (6.66,-0.25);
  % --- H4: particle drawn into the beam
  \fill[black!75] (9.4,-1.55) circle (0.09);
  \draw[->,black!75,thin,densely dashed] (9.3,-1.48) .. controls (8.2,-0.9) .. (6.75,-0.08);
  % --- H6: mechanical drift
  \draw[<->,greydark,thick] (3.2,-1.55) -- (4.4,-1.55);
  % --- annotation boxes
  \tikzset{ann/.style={draw=#1,fill=white,rounded corners=1.5pt,text width=3.55cm,align=left,inner sep=3pt}}
  \node[ann=tealdark,anchor=south west] (h1) at (6.95,2.35)
     {\textbf{\color{tealdark}H1} \Oxy{} reacts with photo-activated surface sites $\to$ absorbing or scattering oxide defects};
  \node[ann=amberdark,anchor=south west] (h2) at (2.95,2.35)
     {\textbf{\color{amberdark}H2} adsorbed organics / siloxanes photo-fixed into a film (oxygen may \emph{protect})};
  \node[ann=reddark,anchor=south west] (h5) at (-1.05,2.35)
     {\textbf{\color{reddark}H5} colour centres in the glass (atmosphere-independent, except \Htwo)};
  \node[ann=cyan!60!blue,anchor=north west] (h3) at (6.95,-2.3)
     {\textbf{\color{cyan!60!blue}H3} adsorbed water assists surface photochemistry or roughening};
  \node[ann=black!75,anchor=north west] (h4) at (2.95,-2.3)
     {\textbf{H4} airborne particles pulled into the high-intensity spot};
  \node[ann=greydark,anchor=north west] (h6) at (-1.05,-2.3)
     {\textbf{\color{greydark}H6} mechanical/thermal drift, epoxy creep (no surface change)};
  \draw[amberdark,thin] (h2.south) -- (6.72,0.27);
  \draw[reddark,thin] (h5.south) -- (5.25,0.09);
  \draw[tealdark,thin] (h1.south) -- (7.55,1.95);
  \draw[cyan!60!blue,thin] (h3.north) -- (7.6,-1.8);
  \draw[black!75,thin] (h4.north) -- (9.35,-1.62);
  \draw[greydark,thin] (h6.north) -- (3.8,-1.62);
\end{tikzpicture}

  \caption{Candidate processes at a UV-illuminated fibre end face (input end shown; the output face sees
  practically the same guided intensity). Only H1--H4 can depend on the surrounding gas. H5 can be
  influenced only by hydrogen, which diffuses into silica; H6 is an artefact. The experiment is designed
  to discriminate between them (\cref{tab:predictions}).}
  \label{fig:mechanisms}
\end{figure}

\section{Methodology and experimental design}\label{sec:methods}

\subsection{Design principles}
Fibre degradation experiments are slow, and many things can change in a lab over several weeks. The design
therefore follows five rules:
\begin{enumerate}[label=\textbf{P\arabic*}]
  \item \textbf{Parallel, not sequential.} All test atmospheres run \emph{at the same time} from
        one light source, in identical optical arms. Laser drifts, lab humidity and seasonal
        effects then affect all arms equally and cancel in the comparison.
  \item \textbf{Change one factor at a time.} The decisive pair of arms differs \emph{only} in
        the presence of \Oxy{}: clean dry synthetic air versus clean dry \Ntwo{}, both passed through the
        same purifier and particle filter (\cref{tab:arms}).
  \item \textbf{Separate the surface from the alignment.} Throughput is recorded before and after
        re-peaking the input coupling, and the back-reflection from the input face is logged
        continuously. Alignment drift (H6) can then be told apart from changes of the end face.
  \item \textbf{Locate the loss.} At the end of every run a fixed sequence of non-destructive
        and destructive steps (\cref{sec:ladder}) splits the loss between input face, output face
        and fibre.
  \item \textbf{Decide in advance.} Metrics, the analysis and the decision thresholds
        (\cref{sec:analysis}) are written down before the confirmatory campaign starts. This
        guards against reading patterns into noisy data after the fact.
\end{enumerate}

\subsection{Phase 0: forensics of fibres already degraded in the lab}\label{sec:forensics}
The cheapest information is already in the lab. In the first two months the student collects
degraded fibres from the group's experiments (ideally 3--6, with as much history as possible:
power, hours of use, wavelength, environment, cleaning and polishing history, PC or APC ends) and:
\begin{enumerate}
  \item images both end faces \emph{before any cleaning} with the fibre inspection microscope (core
        zone, standard magnification and illumination, \cref{app:sops});
  \item measures throughput at low power (\SI{\le 0.3}{mW}) to avoid changing the surfaces further;
  \item runs the cleaning sequence of \cref{sec:ladder} on half of the fibres and keeps the others
        uncleaned for surface analysis (optical and fluorescence microscopy, SEM/EDX, AFM or
        white-light interferometry, Raman, and where available ToF-SIMS or XPS).
\end{enumerate}
In parallel, the student surveys likely contaminant sources near the group's fibre ends: silicone
grease, silicone tubing and adhesives, thermal paste, cable jackets, printed parts and cleaning
practices. Within weeks this shows whether the damage is a \emph{deposit} (and whether it is rich in carbon
or in silicon), a \emph{surface modification} or \emph{nothing visible}. That result already favours some hypotheses
over others and decides which test atmospheres are most informative.

\subsection{Test fibres and end-face preparation}
\begin{itemize}
  \item \textbf{Fibre type.} The same fibre as used in the experiments, so that the results transfer
        directly: Thorlabs P3-375PM-FC patch cables, FC/APC at both ends, built on PM-S350-HP
        pure-silica-core PANDA fibre (\cref{sec:sota-fibres}), ideally from a single batch. Thorlabs
        sells the bare fibre no longer, but a fibre with the same part number and specification is
        available from Coherent (Nufern). Cheaper non-PM SM300 cables, which also have a pure-silica
        core, can serve as extra arms where polarisation does not matter.
  \item \textbf{Format.} \SIrange{1}{2}{m} lengths, connectorised and polished in-house with one
        documented procedure and one epoxy (or bought as identical patch cables). The fibre is coiled
        loosely (bend radius well above the specified minimum) and taped to the table, so that bend
        loss cannot change.
  \item \textbf{Acceptance.} Every end face is inspected and imaged after polishing, and its geometry
        is measured if an end-face interferometer is available. Its initial coupling efficiency is measured
        at low power. Ends that miss the acceptance criteria (no defects in the core zone; coupling
        within a few per cent of the batch median) are re-polished before use.
  \item \textbf{Surface preparation is a variable.} At \SI{405}{nm}, cleaved faces degraded four times
        more slowly than glued-and-polished ones~\cite{Gonschior2013PhD}. Every test end is therefore
        prepared by exactly the same polishing procedure. A comparison of preparation methods (e.g.\ an
        additional final polishing step) is an optional arm in Campaign~II.
  \item \textbf{Traceability.} Each fibre gets an ID. Each end has its own polishing record and image
        set. The assignment of fibres to arms is randomised, and arm positions are swapped between
        runs so that any arm-specific optical problem cannot look like an atmosphere effect.
\end{itemize}

\subsection{Light source and multi-arm testbed}\label{sec:testbed}
\Cref{fig:testbed} shows the testbed. Its main elements are:

\paragraph{Source} A dedicated \SI{375}{nm} single-mode laser diode (\cref{sec:sota-hardware}) with
an optical isolator and anamorphic beam shaping. It frees the group's \SI{369.5}{nm} lasers for
experiments and allows long, uninterrupted runs. The photon energies at \SI{375}{nm} and
\SI{369.5}{nm} differ by only \SI{1.5}{\percent}, and no known silica absorption band lies between them.
Equivalent behaviour is therefore expected, but this assumption is \emph{tested} in Phase~4, which
runs at \SI{369.5}{nm}. If no diode is available, a pick-off of a few tens of milliwatts from an
existing \SI{369.5}{nm} system can feed the testbed.

\paragraph{Power distribution} A cascade of half-wave plates and polarising beam splitters divides the
beam into 4--6 arms with adjustable and equal powers. A further half-wave plate in each arm aligns the
polarisation to the slow axis of the PM fibre. The operational test power is \SI{3}{mW} guided; the
accelerated condition is \SI{10}{mW}. A sentinel arm at \SI{\le 0.1}{mW} in \Ntwo{} measures the
drift of the measurement system itself.

\paragraph{Coupling} An aspheric lens with near-UV anti-reflection coating ($f\approx\SIrange{4.5}{8}{mm}$,
chosen for mode matching, \cref{app:estimates}) focuses the beam through the cell window onto the
fibre end face, at the \ang{3.8} incidence an \ang{8} APC face requires. Commercial launch systems for
FC/APC (e.g.\ Thorlabs FiberPort PAF2A-5A or -7A) or a moulded asphere (e.g.\ C230TMD-A) can be used.
The high-index glasses of such aspheres transmit less well in the near UV, so their transmission at
\SI{369.5}{nm} should be checked. The lens, or the cell, sits on a stable XYZ mount. Automated peak-up with motorised mirror
mounts is a stretch goal; the baseline is a manual re-peak once a day with a fixed procedure.

\paragraph{Detection} A wedged pick-off in each arm feeds a reference photodiode ($P_\text{ref}$) and a
second photodiode that records the light reflected by the input face ($P_\text{back}$; with FC/APC
connectors this detector sits in the specularly reflected beam, which does not retrace the input). The output power $P_\text{out}$
is measured with an integrating-sphere or diffuser-fronted photodiode, which is insensitive to small
changes of the output beam. Silicon photodiodes lose UV responsivity and gain dark current under
prolonged UV exposure~\cite{HamamatsuS1336,Korde1987}. All detectors are therefore operated at low
intensity (attenuated, diffused or shuttered), photodiodes sold for UV resistance are preferred (e.g.\
Hamamatsu S12698 series~\cite{HamamatsuS12698}), and a reference power meter that is otherwise kept
in a drawer cross-checks them weekly. SiC photodiodes, the usual choice for UV-stable detection, are not
suitable here: \SI{369.5}{nm} lies at the edge of their spectral response.

\paragraph{Logging} All photodiodes, laser current and temperature, lab and arm temperatures, humidity,
pressure, flows and \Oxy{} readings are logged at \SI{1}{Hz} (stored as 1-minute means) by a Python
data logger (\cref{app:logging}). Every manual intervention is recorded in an event log with a
timestamp.

\subsection{Test atmospheres}\label{sec:atmospheres}
\begin{table}[t]
\centering\small
\caption{Test atmospheres. Core arms run in Campaign~I; extended arms are selected for Campaign~II
on the basis of Campaign~I and the Phase~0 forensics.}
\label{tab:arms}
\begin{tabularx}{\textwidth}{@{}l L{4.2cm} Y@{}}
\toprule
Arm & Atmosphere at both end faces & Purpose \\
\midrule
\multicolumn{3}{@{}l}{\emph{Core arms}}\\
\arm{AIR}  & ambient lab air, ends open as in the experiments & control; reproduces the problem \\
\arm{CDA}  & synthetic air (\SI{20.9}{\percent} \Oxy{} in \Ntwo{}, zero grade), purified, filtered, dry
           & contains \Oxy{} but no contaminants, particles or water \\
\arm{N2}   & \Ntwo{} (grade 5.0 or LN$_2$ boil-off), same purifier and filter, dry
           & identical to \arm{CDA} except for \Oxy{} $\Rightarrow$ decisive test of H1 vs H2 \\
\arm{ABS}  & sealed printed cell, pre-purged with \Ntwo{}, iron-based \Oxy{} absorber, no flow
           & the low-cost, passive solution to be deployed \\
\midrule
\multicolumn{3}{@{}l}{\emph{Extended arms}}\\
\arm{O2-x} & \Ntwo{} with $x=$ 10, 100, \num{1000}, \num{10000}\,ppm \Oxy{}
           & dose--response, sets the \Oxy{} specification for the cells \\
\arm{WET-N2}  & \Ntwo{} humidified to \SIrange{30}{50}{\percent} RH & water without \Oxy: assists damage (H3)
           or protects against carbon deposits (H2)~\cite{Weiller2010,Parthasarathy2024} \\
\arm{WET-CDA} & synthetic air at lab-like RH, filtered & lab air minus contaminants and particles (H2, H4) \\
\arm{FG}   & forming gas, \SI{5}{\percent} \Htwo{} in \Ntwo{} & hydrogen passivation of defects (H5) \\
\arm{N2+P} & \Ntwo{} cell containing a coupon of the printing material
           & are printed parts a contamination source? (H2) \\
\arm{+SIL} & a silicone coupon in one \arm{CDA} and one \arm{N2} cell
           & oxygen-assisted siloxane deposition (H1)~\cite{Fujifilmpatent2006} \\
\arm{VAC}  & vacuum (optional, if a chamber with fibre feedthrough is free) & link to in-vacuum fibre ends \\
\bottomrule
\end{tabularx}
\end{table}

\Cref{tab:arms} lists the atmospheres. Purged arms receive \SIrange{5}{20}{sccm} per cell through a
common chain of regulator, oxygen and moisture purifier and \SI{3}{nm} particle filter. This is far less
than the \SIrange{0.5}{1}{L/min} that Marciniak et al.\ blew through open collimators to keep particles
out~\cite{Marciniak2017}, but a sealed cell needs no more than a small positive outflow (\cref{app:estimates}). Short runs of
stainless-steel or nylon tubing are used, because polyurethane and fluoropolymer tubing let enough
oxygen through to matter at the ppm level (\cref{app:estimates}). Each outlet passes a humidity sensor, a
VOC sensor and (for \Oxy{}-free arms) an \Oxy{} sensor before reaching an oil bubbler or a long thin
exhaust line that stops back-diffusion. \Oxy{} mixtures for \arm{O2-x} are made by adding synthetic
air to \Ntwo{} with two flow controllers, or are bought as certified mixtures. Humidified arms use a
water bubbler in a split stream. The terms ``getter'' and ``absorber'' are used carefully. Vacuum
non-evaporable getters are not suitable at atmospheric pressure. The cells use iron-based oxygen
absorbers of the kind used in food and museum conservation (\cref{sec:sota-atmosphere}).

\subsection{End-face cells}\label{sec:cells}
Each fibre end that is not in ambient air sits in a small cell (\cref{fig:cell}). The requirements are:
\begin{enumerate}
  \item \textbf{Optical access} through a UV fused-silica window, AR-coated near \SI{370}{nm} and
        slightly wedged and tilted to avoid etalons and back-reflections into the fibre. The window lies
        in the converging beam well away from the focus, so the intensity on it is low. For a
        \SI{1}{mm} window at NA~0.1 the added spherical aberration is negligible
        (\cref{app:estimates}).
  \item \textbf{A small, swept volume} (\SIrange{10}{30}{cm^3}). Gas enters next to the end face and
        leaves downstream, so that anything released by the connector, epoxy or cable is carried
        \emph{away} from the face.
  \item \textbf{Tightness.} The whole connector sits inside the cell and the cable leaves through an
        epoxy-potted gland. Otherwise air diffuses along the jacket to the ferrule. Each cell must pass a
        pressure-decay leak test (target $<\SI{1e-5}{mbar.L/s}$, \cref{app:sops}); at that leak rate a
        typical absorber sachet lasts for months (\cref{app:estimates}).
  \item \textbf{Static variant (\arm{ABS}).} The absorber sits in a pocket behind a porous PTFE
        membrane, out of line of sight of the end face, so that no dust or vapour from the sachet
        reaches it. The cell is first purged with \Ntwo{} to about \SI{1}{\percent} \Oxy{}. The absorber
        then only has to remove a small residue, and the pressure in the sealed cell drops by about
        \SI{10}{mbar} instead of \SI{210}{mbar}, which would otherwise draw air in through any leak.
  \item \textbf{Verification.} An optical \Oxy{} sensor spot (read through a window, so that no
        electronics sit inside the cell) and a colour \Oxy{} indicator monitor the atmosphere.
  \item \textbf{Clean materials.} Printed parts are a possible source of organic vapour (H2). Prototypes
        are printed in PETG or ASA with \SI{100}{\percent} infill and several perimeters, sealed with a
        low-outgassing epoxy and baked under \Ntwo{} purge. The \arm{N2+P} arm tests directly whether the
        material harms the end face. O-rings are fitted dry or with a minimal film of PFPE grease, and
        silicone grease is banned because of siloxane contamination. The gland is potted with a
        low-outgassing epoxy.
\end{enumerate}
The student designs the cells in CAD with the dimensions of the group's fibre mounts. Two designs are
worth prototyping early: (a) a cap that encloses the connector and screws onto a standard
fibre-launch receptacle, and (b) a box that encloses the whole launch (lens included). Design (a) has a
smaller volume; design (b) is easier to seal.

\subsection{Diagnostics}
\begin{table}[t]
\centering\small
\caption{Measurements, their purpose and cadence.}
\label{tab:diagnostics}
\begin{tabularx}{\textwidth}{@{}L{4.1cm} Y L{2.6cm}@{}}
\toprule
Measurement & What it tells us & When \\
\midrule
$\eta = P_\text{out}/P_\text{ref}$ (normalised) & total throughput change of each arm & continuous, \SI{1}{Hz}\\
re-peak gain $\eta_\text{after}/\eta_\text{before}$ & alignment drift (H6) vs end-face change & daily or automated \\
$P_\text{back}/P_\text{ref}$ & optical change of the input face (film, roughness) & continuous \\
\Oxy, RH, VOC index, $T$, $p$, flow & that each atmosphere is what it should be & continuous \\
output beam profile and power through a mode-selective aperture & loss of beam quality (rises in $M^2$ have
been seen at \SI{405}{nm}~\cite{Gonschior2013PhD}) that a large detector does not register & every \SI{30}{min} \\
visible luminescence from the input face (spectrometer, optional) & fluorescence of growing organic films; defect PL bands & every \SI{30}{min} \\
end-face microscopy (bright field, fluorescence) & deposits, particles, damage & start and end of each run \\
cleaning sequence (\cref{sec:ladder}) & splits the loss between faces and bulk & end of each run \\
SEM/EDX, AFM or interferometry, Raman, ToF-SIMS/XPS & composition and thickness of films & selected ends \\
\bottomrule
\end{tabularx}
\end{table}
\Cref{tab:diagnostics} summarises the measurements. Two of them are cheap but very useful. First, the
\emph{power reflected by the input face} is \SI{3.7}{\percent} of the incident power. With the
\ang{8} FC/APC faces of the group's cables this reflection does not retrace the input beam. It leaves
\ang{\approx 24} away from it (\cref{app:estimates}) and can be caught with a small mirror or
photodiode. Even a thin film changes it through thin-film interference and scattering, so it is a
sensitive in-situ monitor of the input face. Second, \emph{organic films fluoresce} under near-UV excitation, while clean silica
hardly does. A long-pass filter and a small spectrometer that looks back through the coupling lens can
therefore see a film start to grow long before the throughput drops noticeably.

\subsubsection{End-of-run cleaning sequence}\label{sec:ladder}
At the end of each run the transmission is measured at low power after each of the following steps,
always with a re-peak and first at the output end, then at the input end:
\begin{enumerate}
  \item no-touch inspection and imaging (preserves evidence);
  \item dry clean (filtered gas jet, then a dry cleaning tool);
  \item solvent clean (reagent-grade isopropanol, lint-free wipe; acetone as a second solvent);
  \item oxidative clean (\Oxy{} plasma or UV/ozone, if available), which removes organic but not
        inorganic (e.g.\ siloxane-derived \ce{SiO_x}) films;
  \item re-polish, recording the removed length.
\end{enumerate}
The loss recovered at each step assigns the damage to a face and to a category (loose particles,
soluble film, oxidisable film, bonded or structural damage). Any loss left after both ends have been
re-polished is attributed to the fibre itself (H5). Where an arm has two fibres, one goes through
this sequence and the other is kept uncleaned for surface analysis.

\subsection{Experimental campaigns}\label{sec:campaigns}
\begin{description}[style=nextline,leftmargin=1.2em,itemsep=4pt]
\item[Phase 1 -- baseline in ambient air (months 4--5)] Three to four \arm{AIR} arms at \SI{1}{mW},
  \SI{3}{mW} and \SI{10}{mW}, plus the sentinel. Outputs: (i) proof that the testbed reproduces the
  degradation seen in the lab; (ii) the power law $r\propto P^{n}$ of the initial degradation rate. An
  exponent $n\approx2$ points to a two-photon process; (iii) fibre-to-fibre scatter; (iv) the run length
  needed in the later campaigns. Raising the power only accelerates the test if the mechanism stays the
  same. Surface restructuring at \SI{405}{nm} was reported above \SI{\sim0.5}{MW/cm^2}~\cite{Gonschior2013JNCS},
  close to the \SI{0.44}{MW/cm^2} reached here at \SI{10}{mW}. The end faces from the \SI{3}{mW} and
  \SI{10}{mW} arms are therefore compared by microscopy before \SI{10}{mW} is used as an accelerated condition.
\item[Phase 2 -- Campaign I, atmosphere screening (months 6--7)] \arm{AIR}, \arm{CDA}, \arm{N2} and
  \arm{ABS} in parallel at the accelerated power (or \SI{3}{mW} if Phase~1 shows that \SI{10}{mW} changes
  the mechanism, e.g.\ thermal runaway). Each run lasts until \arm{AIR} has lost \SI{\sim30}{\percent} or
  for three weeks, whichever comes first. If time allows, a second run with permuted arm positions follows.
\item[Phase 3 -- Campaign II, confirmation and mechanism (months 8--9)] Repeat the key contrast at
  \SI{3}{mW} with at least two fibres per condition, and add the most informative extended arms
  (\cref{tab:arms}): the \Oxy{} dose--response if \Oxy{} matters, \arm{WET-N2}/\arm{WET-CDA} and \arm{N2+P}
  if it does not. A crossed configuration (input in \Ntwo{}, output in air, and vice versa) tests whether
  both faces respond in the same way. Marciniak et al.\ observed that the purged and unpurged ends of
  the same fibre behaved differently~\cite{Marciniak2017}.
\item[Phase 4 -- demonstration at 369.5\,nm (months 9--11)] The best cell design is fitted to both ends
  of a working \SI{369.5}{nm} fibre in one of the group's experiments, with a control fibre or the fibre's
  own history as reference, for at least eight weeks. This phase also checks the transfer from
  \SI{375}{nm} to \SI{369.5}{nm} and whether the cell is practical in use (alignment stability, access,
  window losses).
\end{description}

\begin{decisionbox}[How the outcome of Campaign I steers the project]
\begin{itemize}[leftmargin=1em]
  \item \textbf{\arm{N2} and \arm{ABS} stable, \arm{CDA} degrades like \arm{AIR}}: oxygen takes part
        (H1). Check the deposit chemistry: if it is rich in Si and O, siloxanes are involved, and their
        sources near the fibres should be removed anyway. Then measure the \Oxy{} dose--response,
        specify the cell atmosphere and deploy \arm{ABS} cells. This is the outcome the group hopes for.
  \item \textbf{\arm{CDA} and \arm{N2} both stable, \arm{AIR} degrades}: something in lab air other than
        \Oxy{} is responsible: particles, organic vapours or water (H4, H2, or a water-assisted H3).
        Campaign~II separates them. A sealed or purged \emph{clean} cell solves the practical problem, with
        or without oxygen. This is what the published nitrogen-purge results would suggest~\cite{Marciniak2017}.
  \item \textbf{\arm{CDA} stable, \arm{N2} or \arm{ABS} degrade}: carbonaceous contamination with
        oxygen as a protective agent (H2). This is the outcome that space-laser, UV-cavity and laser-diode
        experience predicts (\cref{sec:sota-surface}). Oxygen absorbers would be counter-productive; the
        recommendation becomes a clean, dry (or humidified) air purge.
  \item \textbf{No atmosphere helps}: restructuring of the silica, colour centres or an artefact (H3, H5,
        H6). Shift the effort to end-face and fibre engineering. Coreless end caps lower the intensity at
        the glass--gas interface by two to three orders of magnitude (\cref{app:estimates}). Other options are
        hydrogen-loaded and UV-cured fibre and hollow-core fibre (\cref{sec:sota-fibres}).
\end{itemize}
Every branch gives a publishable answer and a concrete recommendation for the group.
\end{decisionbox}

\subsection{Data analysis and decision rules}\label{sec:analysis}
\paragraph{Metrics} The primary quantity is the normalised, re-peaked throughput
$\eta(t)=[P_\text{out}/P_\text{ref}](t)/[P_\text{out}/P_\text{ref}](t_0)$, corrected for detector
temperature coefficients, with flagged intervals (interventions, outages) excluded. The \emph{primary
endpoint} is the initial degradation rate $r=-\mathrm{d}\eta/\mathrm{d}t$ in \si{\percent} per
\SI{1000}{h} at \SI{3}{mW}. Secondary endpoints are the time to \SI{10}{\percent} loss, the fit
parameters of a phenomenological model (e.g.\ a stretched exponential
$\eta=1-A\{1-\exp[-(t/\tau)^{\beta}]\}$ in time or dose), the change of $P_\text{back}$, and the outcome of
the cleaning sequence.

\paragraph{Statistics} Confidence intervals come from a block bootstrap, which allows for
autocorrelated data. Conditions are compared through rate ratios. With replicate fibres, a
mixed-effects model with fibre as random effect separates the atmosphere effect from fibre-to-fibre
scatter. Phase~1 provides the scatter needed to choose the number of replicates.

\paragraph{Decision rules (fixed before Campaign II)}
\begin{itemize}
  \item \textbf{\Oxy{} effect established:} $r_\text{N2}/r_\text{CDA}$ differs from 1 by at least a factor of
        two, with the 95\,\% confidence interval excluding 1, in two independent runs.
  \item \textbf{Degradation halted:} the 95\,\% upper limit on $|r|$ is below \SI{1}{\percent} per
        \SI{1000}{h} at \SI{3}{mW} (equivalent to $<\SI{10}{\percent}$ loss per year of continuous use),
        at least twenty times below \arm{AIR}, with no visible change of either end face. With a daily
        scatter of \SI{0.3}{\percent} this limit can be reached in three to six weeks (\cref{app:estimates}).
\end{itemize}

\subsection{Success criteria and deliverables}
\begin{enumerate}[label=\textbf{S\arabic*}]
  \item The location and physical nature of the degradation are established for the group's fibres (RQ1).
  \item The effect of \Oxy{} on the degradation rate is measured with a stated confidence interval, or
        excluded at the level above (RQ2), and the roles of water, contaminants and particles are tested (RQ3).
  \item A cell design is qualified: \Oxy{} $<\SI{100}{ppm}$ (target $<\SI{10}{ppm}$) and RH
        $<\SI{5}{\percent}$ for at least eight weeks without maintenance. It is tested on a working
        \SI{369.5}{nm} link (RQ4). If the atmosphere helps, the degradation-halted criterion is met.
  \item Open deliverables: CAD files and bill of materials of the cells, standard operating procedures, the
        data set with analysis code, the MSc thesis, and (if the results allow) a short paper.
\end{enumerate}


\begin{figure}[p]
  \centering
  \resizebox{\textwidth}{!}{% Multi-arm accelerated-ageing testbed: one source, N identical arms, one atmosphere per arm.
\begin{tikzpicture}[font=\sffamily\scriptsize,>=Stealth,x=1cm,y=1cm,
    beam/.style={uvviolet,line width=0.9pt},
    comp/.style={draw=greydark,fill=white,rounded corners=1pt,inner sep=2pt,align=center},
    cell/.style={draw=tealdark,fill=teallight,rounded corners=2pt,thick},
    pd/.style={draw=greydark,fill=gray!35},
    gas/.style={tealdark,line width=0.7pt},
    sig/.style={greydark!80,densely dashed,line width=0.4pt}]
  % ------------------------------------------------------------ source and splitter column
  \node[comp,fill=uvlight,minimum width=2.1cm,minimum height=0.75cm] (laser) at (1.2,8.25)
       {\textbf{laser}\\375\,nm diode or 369.5\,nm};
  \node[comp,minimum width=2.1cm] (shape) at (1.2,7.15) {beam shaping (+ isolator)};
  \draw[beam,->] (laser.south) -- (shape.north);
  \draw[beam] (shape.south) -- (1.2,1.5);
  \foreach \y/\lab/\desc/\celltype in {6.0/1/{ambient lab air (open)}/0,
                                       4.5/2/{clean dry synthetic air, purged}/1,
                                       3.0/3/{clean dry \Ntwo{}, purged}/1,
                                       1.5/4/{sealed cell + \Oxy{} absorber}/2}{
    % splitter: lambda/2 plate + PBS (last arm: mirror)
    \ifnum\lab<4
      \draw[greydark,fill=uvlight] (1.2-0.22,\y-0.22) rectangle (1.2+0.22,\y+0.22);
      \draw[greydark] (1.2-0.22,\y-0.22) -- (1.2+0.22,\y+0.22);
      \draw[greydark,fill=white] (0.95,\y+0.42) rectangle (1.45,\y+0.48);
    \else
      \draw[greydark,line width=1.6pt] (1.2-0.2,\y-0.2) -- (1.2+0.2,\y+0.2);
    \fi
    \draw[beam,->] (1.42,\y) -- (2.75,\y);
    % pick-off wedge with reference / back-reflection photodiodes
    \draw[greydark,fill=cyan!10] (2.85,\y-0.32) -- (3.0,\y-0.32) -- (3.1,\y+0.32) -- (2.95,\y+0.32) -- cycle;
    \draw[beam,densely dotted] (2.98,\y) -- (2.98,\y+0.55);
    \fill[pd] (2.80,\y+0.55) rectangle (3.16,\y+0.66);
    \draw[beam,densely dotted] (3.0,\y) -- (3.0,\y-0.55);
    \fill[pd] (2.82,\y-0.66) rectangle (3.18,\y-0.55);
    \draw[beam,->] (3.1,\y) -- (3.95,\y);
    % focusing asphere
    \draw[greydark,fill=cyan!10] (4.15,\y) ellipse (0.08 and 0.33);
    % input-end atmosphere
    \ifnum\celltype=0
      \draw[greydark!60,densely dashed,rounded corners=2pt] (4.55,\y-0.42) rectangle (5.95,\y+0.42);
    \else
      \draw[cell] (4.55,\y-0.42) rectangle (5.95,\y+0.42);
      \fill[cyan!25] (4.62,\y-0.3) rectangle (4.72,\y+0.3);
    \fi
    \fill[uvviolet!25] (4.23,\y+0.22) -- (5.5,\y+0.03) -- (5.5,\y-0.03) -- (4.23,\y-0.22) -- cycle;
    \fill[gray!55] (5.5,\y-0.09) rectangle (5.95,\y+0.09);
    % fibre (coil) between the two cells
    \draw[uvviolet!70!black,line width=0.8pt] (5.95,\y) -- (6.45,\y)
      .. controls (6.75,\y+0.38) and (7.15,\y+0.38) .. (7.3,\y)
      .. controls (7.45,\y-0.38) and (7.85,\y-0.38) .. (8.0,\y)
      .. controls (8.15,\y+0.38) and (8.55,\y+0.38) .. (8.7,\y) -- (9.2,\y);
    % output-end atmosphere
    \ifnum\celltype=0
      \draw[greydark!60,densely dashed,rounded corners=2pt] (9.2,\y-0.42) rectangle (10.6,\y+0.42);
    \else
      \draw[cell] (9.2,\y-0.42) rectangle (10.6,\y+0.42);
      \fill[cyan!25] (10.43,\y-0.3) rectangle (10.53,\y+0.3);
    \fi
    \fill[gray!55] (9.2,\y-0.09) rectangle (9.65,\y+0.09);
    \fill[uvviolet!25] (9.65,\y+0.03) -- (11.25,\y+0.3) -- (11.25,\y-0.3) -- (9.65,\y-0.03) -- cycle;
    % output detector (integrating sphere)
    \draw[pd] (11.6,\y) circle (0.3);
    \fill[white] (11.3,\y-0.07) rectangle (11.36,\y+0.07);
    % label
    \node[anchor=west,align=left] at (12.05,\y) {\textbf{arm \lab}\\\desc};
    \ifnum\celltype=2
      \fill[amberdark!75] (5.15,\y-0.40) rectangle (5.45,\y-0.28);
      \fill[orange] (5.25,\y+0.33) circle (0.05);
      \fill[amberdark!75] (9.85,\y-0.40) rectangle (10.15,\y-0.28);
      \fill[orange] (9.95,\y+0.33) circle (0.05);
    \fi
  }
  % ------------------------------------------------------------ annotations of the first arm
  \node[anchor=south,align=center] at (4.15,6.45) {asphere\\$f\!\approx\!5$\,mm};
  \node[anchor=south,align=center] at (5.25,6.45) {input\\end};
  \node[anchor=south] at (7.6,6.45) {1--2\,m PM fibre};
  \node[anchor=south,align=center] at (9.9,6.45) {output\\end};
  \node[anchor=south,align=center] at (11.6,6.45) {$P_\text{out}$};
  \node[anchor=north,align=center] at (3.0,0.78) {pick-off wedge:\\$P_\text{ref}$ (up), $P_\text{back}$ (down)};
  % purged cells: gas in / out stubs
  \foreach \y in {4.5,3.0}{
    \draw[gas,->] (5.05,\y+0.68) -- (5.05,\y+0.42);
    \draw[gas,->] (5.45,\y+0.42) -- (5.45,\y+0.68);
    \draw[gas,->] (9.7,\y+0.68) -- (9.7,\y+0.42);
    \draw[gas,->] (10.1,\y+0.42) -- (10.1,\y+0.68);
  }
  % ------------------------------------------------------------ gas panel
  \node[comp,fill=teallight,draw=tealdark,minimum width=2.4cm,align=left] (cyl) at (1.75,-0.95)
       {\Ntwo{} (5.0 or LN$_2$ boil-off)\\synthetic air (zero grade)\\opt.: forming gas, \Oxy/\Ntwo{} mix};
  \node[comp,draw=tealdark,align=center] (pur) at (5.0,-0.95) {regulators, purifier,\\3\,nm particle filter};
  \node[comp,draw=tealdark,align=center] (flow) at (7.85,-0.95) {flow control\\5--20\,sccm per cell};
  \draw[gas,->] (cyl.east) -- (pur.west);
  \draw[gas,->] (pur.east) -- (flow.west);
  \draw[gas,->] (flow.east) -- (9.35,-0.95) node[anchor=west,align=left,tealdark] {to the cells\\of arms 2 and 3};
  \node[tealdark,anchor=north west,align=left] at (0.55,-1.65)
       {cell outlets $\to$ RH sensor, trace \Oxy{} analyser, VOC sampling $\to$ long exhaust line or check valve (no oil bubbler)};
  % ------------------------------------------------------------ data acquisition
  \node[comp,minimum width=2.4cm,align=left,anchor=west] (daq) at (12.0,-0.95)
       {\textbf{DAQ + PC}\\photodiodes, \Oxy, RH, $T$, $p$,\\flows, laser current};
  % ------------------------------------------------------------ legend
  \begin{scope}[shift={(0.55,-2.45)}]
    \draw[cell] (0,0) rectangle (0.35,0.22); \node[anchor=west] at (0.4,0.11) {end-face cell (3D-printed, UV window)};
    \fill[amberdark!75] (5.1,0.04) rectangle (5.4,0.18); \node[anchor=west] at (5.45,0.11) {\Oxy{} absorber};
    \fill[orange] (7.5,0.11) circle (0.05); \node[anchor=west] at (7.6,0.11) {\Oxy{} sensor spot};
    \fill[pd] (9.75,0.04) rectangle (10.05,0.18); \node[anchor=west] at (10.1,0.11) {photodiode};
    \draw[gas] (11.7,0.11) -- (12.1,0.11); \node[anchor=west] at (12.15,0.11) {clean gas};
  \end{scope}
\end{tikzpicture}
}
  \caption{Multi-arm accelerated-ageing testbed. One near-UV source feeds identical arms that differ only
  in the atmosphere at the two end faces. Each arm records the reference, back-reflected and transmitted
  powers. Purged arms share one gas-purification chain; the static arm uses an oxygen absorber and an
  optical \Oxy{} sensor spot. Arm 1 has no cells and reproduces the present lab conditions.}
  \label{fig:testbed}
\end{figure}

\begin{figure}[t]
  \centering
  % Concept cross-section of a 3D-printed end-face atmosphere cell (input end shown).
\begin{tikzpicture}[font=\sffamily\scriptsize,>=Stealth,x=1cm,y=1cm,
    lab/.style={align=left,inner sep=1.5pt},
    lead/.style={greydark,thin}]
  % ------------------------------------------------------------ printed body
  \fill[pattern=north east lines,pattern color=gray!55] (0,0) rectangle (9,4);
  \draw[greydark,thick] (0,0) rectangle (9,4);
  \fill[white] (0.45,0.45) rectangle (8.55,3.55);
  \draw[greydark] (0.45,0.45) rectangle (8.55,3.55);
  % ------------------------------------------------------------ absorber pocket (below cavity)
  \fill[pattern=north east lines,pattern color=gray!55] (3.7,-1.1) rectangle (6.5,0);
  \draw[greydark,thick] (3.7,-1.1) rectangle (6.5,0);
  \fill[white] (4.05,-0.8) rectangle (6.15,0.45);
  \draw[greydark] (4.05,-0.8) -- (4.05,0.45) (6.15,-0.8) -- (6.15,0.45) (4.05,-0.8) -- (6.15,-0.8);
  \draw[tealdark,very thick,densely dashed] (4.05,0.45) -- (6.15,0.45);
  \fill[amberdark!75,rounded corners=2pt] (4.35,-0.65) rectangle (5.85,-0.05);
  \node[white] at (5.1,-0.35) {\textbf{absorber}};
  \fill[black] (3.85,-1.1) circle (0.06) (6.35,-1.1) circle (0.06);
  % ------------------------------------------------------------ front window + O-ring
  \fill[white] (0,1.35) rectangle (0.45,2.65);
  \fill[cyan!25] (0.12,1.25) rectangle (0.36,2.75);
  \draw[cyan!50!gray] (0.12,1.25) rectangle (0.36,2.75);
  \fill[black] (0.43,1.33) circle (0.07) (0.43,2.67) circle (0.07);
  % ------------------------------------------------------------ viewing window + O2 sensor spot
  \fill[white] (4.6,3.55) rectangle (5.6,4);
  \fill[cyan!25] (4.5,3.62) rectangle (5.7,3.86);
  \draw[cyan!50!gray] (4.5,3.62) rectangle (5.7,3.86);
  \fill[orange] (4.85,3.56) rectangle (5.05,3.62);
  \fill[green!50!black] (5.2,3.56) rectangle (5.4,3.62);
  \draw[greydark,thick] (4.95,3.86) -- (4.95,4.55);
  \node[draw=greydark,fill=white,rounded corners=1pt,lab] at (4.95,4.85) {optical \Oxy{} reader (outside)};
  % ------------------------------------------------------------ lens and beam
  \fill[uvviolet!20] (-2.8,1.45) -- (-1.25,1.45) -- (2.35,1.97) -- (2.35,2.03) -- (-1.25,2.55) -- (-2.8,2.55) -- cycle;
  \draw[uvviolet!70] (-2.8,1.45) -- (-1.25,1.45) -- (2.35,1.97) (-2.8,2.55) -- (-1.25,2.55) -- (2.35,2.03);
  \draw[greydark,fill=cyan!10] (-1.25,2.0) ellipse (0.13 and 0.75);
  \draw[->,uvviolet,very thick] (-2.7,2.0) -- (-1.75,2.0);
  % ------------------------------------------------------------ receptacle, ferrule, connector, cable
  \fill[gray!70] (2.75,1.5) rectangle (3.35,2.5);
  \fill[pattern=north east lines,pattern color=gray!55] (2.85,0.45) rectangle (3.25,1.5);
  \draw[greydark] (2.85,0.45) rectangle (3.25,1.5);
  \fill[gray!25] (2.35,1.9) rectangle (3.6,2.1);
  \draw[gray!60!black] (2.35,1.9) rectangle (3.6,2.1);
  \fill[uvviolet] (2.35,2.0) circle (0.055);
  \fill[gray!55!black] (3.35,1.62) rectangle (5.3,2.38);
  \fill[gray!45!black] (5.3,1.7) -- (6.5,1.88) -- (6.5,2.12) -- (5.3,2.3) -- cycle;
  \draw[black!70,line width=1.6pt] (6.5,2.0) -- (10.3,2.0);
  \fill[gray!80!black] (8.5,1.7) rectangle (9.15,2.3);
  \fill[amberdark!40] (8.62,1.85) rectangle (9.03,2.15);
  % ------------------------------------------------------------ gas ports
  \fill[white] (1.0,3.55) rectangle (1.4,4.0);
  \draw[greydark] (1.0,3.55) -- (1.0,4.35) (1.4,3.55) -- (1.4,4.35);
  \draw[->,tealdark,thick] (1.2,4.55) -- (1.2,3.25) .. controls (1.25,2.7) .. (2.15,2.25);
  \fill[white] (7.6,0) rectangle (8.0,0.45);
  \draw[greydark] (7.6,0.45) -- (7.6,-0.35) (8.0,0.45) -- (8.0,-0.35);
  \draw[->,tealdark,thick] (6.9,1.2) .. controls (7.7,0.9) .. (7.8,0.35) -- (7.8,-0.55);
  % ------------------------------------------------------------ labels
  \node[lab,anchor=south west] (lwin) at (-2.85,3.25) {UV fused-silica window,\\AR-coated at 370\,nm, wedged;\\O-ring seal};
  \draw[lead] (lwin.south east) ++(-0.6,0) -- (0.24,2.75);
  \node[lab,anchor=north] at (-1.25,1.15) {lens outside:\\$f\!\approx\!5$--8\,mm};
  \node[lab,anchor=south west,tealdark] at (1.35,4.2) {\textbf{gas in} (sweeps the end face first)};
  \node[lab,anchor=north west,tealdark] at (8.1,-0.15) {\textbf{gas out}\\(downstream)};
  \node[lab,anchor=north east] (lfer) at (2.55,1.0) {end face under test};
  \draw[lead] (lfer.north east) ++(-0.1,0) -- (2.33,1.93);
  \node[lab,anchor=north west] (lrec) at (1.55,-0.25) {FC receptacle\\on internal post};
  \draw[lead] (2.75,-0.25) -- (3.0,0.6);
  \node[lab,anchor=north west] (lwall) at (-0.95,-0.25) {printed wall\\(sealed, baked)};
  \draw[lead] (0.2,-0.25) -- (0.25,0.25);
  \node[lab,anchor=north west] at (6.7,-1.0) {absorber (+ desiccant) behind\\porous PTFE membrane};
  \draw[lead] (6.7,-1.05) -- (6.2,-0.3);
  \node[lab,anchor=south west] at (6.75,3.95) {\Oxy{} sensor spot and\\indicator behind a window};
  \draw[lead] (6.75,4.15) -- (5.4,3.6);
  \node[lab,anchor=south west] at (7.15,2.55) {cable gland,\\epoxy-potted};
  \draw[lead] (8.3,2.6) -- (8.75,2.3);
  \node[lab,anchor=north west] at (3.55,3.35) {free volume 10--30\,cm$^3$};
\end{tikzpicture}

  \caption{Concept of a 3D-printed end-face cell (cross-section, input end). The whole connector is
  enclosed and the cable is sealed in a potted gland, because air would otherwise diffuse along the cable
  to the ferrule. Gas enters next to the end face and leaves downstream. In the static variant the inlet
  and outlet are closed after an \Ntwo{} pre-purge and the absorber in the pocket removes the residual
  oxygen. The window sits in the converging beam, far enough from the focus that the intensity on it is
  low.}
  \label{fig:cell}
\end{figure}

\section{Work plan}\label{sec:workplan}

The plan assumes a full-time, 12-month project (\cref{fig:gantt}). Items that need ordering are listed
in \cref{sec:resources} and should be ordered in month~1, because optics, gas fittings and sensors
can have lead times of several weeks. A compressed 6-month version is given at the end of this section.

\begin{figure}[t]
  \centering
  % 12-month work plan in half-month units (1 unit = ~2 weeks).
\begin{ganttchart}[
    x unit=0.43cm, y unit chart=0.5cm, y unit title=0.5cm,
    hgrid={*1{draw=gray!20}}, vgrid={*1{draw=gray!12},*1{draw=gray!35}},
    time slot format=simple, canvas/.append style={draw=none},
    title/.append style={fill=uvviolet!12,draw=gray!40}, title label font=\sffamily\scriptsize\bfseries,
    bar/.append style={fill=uvviolet!55,draw=uvviolet!80!black}, bar height=0.55,
    bar label font=\sffamily\scriptsize, bar label node/.append style={align=right},
    group/.append style={fill=tealdark,draw=tealdark}, group label font=\sffamily\scriptsize\bfseries,
    group right shift=0, group top shift=0.35, group height=0.3, group peaks height=0.15,
    milestone/.append style={fill=amberdark,draw=amberdark!70!black}, milestone height=0.55,
    milestone label font=\sffamily\scriptsize\itshape\color{amberdark!80!black},
    milestone label node/.append style={align=right},
    link/.style={-Stealth,draw=gray!60}
  ]{1}{24}
  \gantttitle{Month}{24} \\
  \gantttitlelist{1,...,12}{2} \\
  \ganttgroup{Set-up}{1}{10} \\
  \ganttbar{WP0 Induction, safety, reading}{1}{3} \\
  \ganttbar{WP1 Forensics of used fibres}{2}{4} \\
  \ganttbar{WP2 Testbed + data logging}{3}{8} \\
  \ganttbar{WP3 End-face cells: design, print, qualify}{4}{10} \\
  \ganttmilestone{MS1 Plan + risk assessments signed off}{2} \\
  \ganttmilestone{MS2 Forensic report; MS3 testbed qualified}{8} \\
  \ganttgroup{Experiments}{7}{22} \\
  \ganttbar{WP4 Baseline in air: power + dose law}{7}{10} \\
  \ganttbar{WP5 Campaign I: atmosphere screening}{11}{14} \\
  \ganttbar{WP6 Campaign II: confirmation + dose-response}{15}{18} \\
  \ganttbar{WP7 Ex-situ surface analysis}{12}{20} \\
  \ganttbar{WP8 Demonstration on the 369.5\,nm system}{17}{22} \\
  \ganttmilestone{MS4 Cells qualified ($<\!100$\,ppm \Oxy{} for 2 weeks)}{10} \\
  \ganttmilestone{MS5 Campaign I decision}{14} \\
  \ganttmilestone{MS6 Campaign II result}{18} \\
  \ganttgroup{Write-up}{8}{24} \\
  \ganttbar[bar/.append style={fill=uvviolet!25}]{WP9 Methods chapter (while building)}{8}{12} \\
  \ganttbar{WP9 Analysis, thesis, paper draft}{18}{24} \\
  \ganttmilestone{MS7 Thesis submitted}{24}
\end{ganttchart}

  \caption{Work packages (WP) and milestones (MS) for a 12-month project in half-month resolution. The
  methods chapter is drafted while the testbed is being built; the analysis and thesis writing overlap
  with the demonstration phase.}
  \label{fig:gantt}
\end{figure}

\begin{table}[t]
\centering\small
\caption{Work packages, main outputs and milestones.}
\label{tab:wps}
\begin{tabularx}{\textwidth}{@{}l Y L{5.0cm}@{}}
\toprule
WP & Content & Output / milestone \\
\midrule
WP0 & inductions (laser, gas, chemical, workshop), literature, project plan & MS1: plan, risk assessments
      and order list agreed (end of month~1) \\
WP1 & forensics of degraded fibres from the lab (\cref{sec:forensics}) & MS2: short forensic report with
      images; choice of extended arms \\
WP2 & testbed: source, splitting, 4--6 arms, detectors, logger, sentinel arm & MS3: testbed qualified
      (drift of sentinel $<\SI{0.3}{\percent}$/day; logging robust for one week) \\
WP3 & end-face cells: CAD, printing, sealing, leak and \Oxy{} hold tests, material screening
    & MS4: cells qualified ($<\SI{100}{ppm}$ \Oxy{} for two weeks, leak rate $<\SI{1e-5}{mbar.L/s}$) \\
WP4 & Phase 1 baseline in air: power and dose dependence, fibre scatter & degradation law; run length
      for later campaigns \\
WP5 & Phase 2 / Campaign I: \arm{AIR}, \arm{CDA}, \arm{N2}, \arm{ABS} & MS5: decision on Campaign~II
      (decision box in \cref{sec:campaigns}) \\
WP6 & Phase 3 / Campaign II: replicates, extended arms, crossed input/output & MS6: confirmatory result
      against pre-registered criteria \\
WP7 & ex-situ surface analysis of selected ends & identification of films or damage \\
WP8 & Phase 4: cells on a working \SI{369.5}{nm} link for $\geq 8$ weeks & demonstration result,
      practical guidelines \\
WP9 & methods chapter (months 4--6), analysis, thesis, paper draft & MS7: thesis submitted \\
\bottomrule
\end{tabularx}
\end{table}

\subsection{Decision points}
\begin{description}[style=nextline,leftmargin=1.2em,itemsep=3pt]
  \item[D1, end of WP4: is the degradation fast enough to study?] If \arm{AIR} at \SI{3}{mW} loses less
        than about \SI{5}{\percent} in two weeks, the campaigns run at \SI{10}{mW}. If necessary the
        intensity is raised further with a smaller mode-field fibre or by running two arms in series on
        one source. If the testbed does not reproduce the lab degradation at all, the difference
        between testbed and lab (wavelength, contaminant sources near the experiments, cleaning
        practice) becomes a research question in itself, and the forensic results guide the next steps.
  \item[D2, end of WP5: which mechanism?] Follow the branches of the decision box in \cref{sec:campaigns}.
  \item[D3, end of WP6: what is deployed?] The cell type (purged or static, with or without oxygen) for
        the Phase~4 demonstration follows from the confirmatory result.
\end{description}

\subsection{Compressed 6-month variant}
If the project has to fit into six months, keep the logic but drop breadth:
month~1: inductions, forensics, ordering; month~2: two to four arms and simple purged cells
(no static cell yet); month~3: short baseline at \SI{10}{mW}; months~4--5: one campaign with \arm{AIR},
\arm{CDA}, \arm{N2} (plus \arm{ABS} if ready); month~6: cleaning sequence, analysis, writing. The
demonstration on a working link and the extended arms become recommendations for a follow-up project.

\begin{keybox}[Minimum viable thesis]
Even if the full testbed is delayed, the following is enough for a good thesis: the forensic analysis
of used fibres, plus one well-controlled parallel comparison of ambient air, clean dry air and dry
\Ntwo{} at the end faces, with the cleaning sequence at the end. The extended arms, the static absorber
cells and the demonstration are the parts that turn it into an excellent thesis and a paper.
\end{keybox}

\section{Resources and indicative budget}\label{sec:resources}

\subsection{What the group is likely to have already}
Check these in month~1: fibre polishing and inspection equipment; spare \SI{369.5}{nm} power on one of the
laser systems (for Phase~4 and as a fall-back source); optical mounts, mirrors, wave plates and beam
splitters for the near UV; power meters; the laboratory \Ntwo{} supply; 3D printers (FDM, possibly SLA)
in the group or the workshop; a fusion splicer (for end caps, optional); and access to SEM/EDX, AFM or
white-light interferometry, Raman and XPS/ToF-SIMS through university facilities.

\subsection{Items to buy}
\Cref{tab:budget} lists the main items. Thorlabs prices are US list prices read on 18~October 2026
(UK prices and university discounts differ). Items marked $\approx$ are rough estimates that should be
replaced by quotes.

\begin{longtable}{@{}L{4.4cm} L{6.6cm} r@{}}
\caption{Main items for the testbed and cells (US dollars).}\label{tab:budget}\\
\toprule
Item & Example and quantity & Cost \\
\midrule
\endfirsthead
\toprule
Item & Example and quantity & Cost \\
\midrule
\endhead
\bottomrule
\endfoot
\multicolumn{3}{@{}l}{\emph{Light source and distribution}}\\
\SI{375}{nm} single-mode diode & Thorlabs L375P70MLD (Nichia NDU4116, \SI{70}{mW})~\cite{ThorlabsL375} & 5\,589 \\
TEC mount & Thorlabs LDM56F (supports the F pin code) & 810 \\
diode current and TEC controller & one combined unit & $\approx$\,2\,000--3\,000 \\
splitting and steering optics & half-wave plates, PBS cubes, mirrors, mounts for 4--6 arms & $\approx$\,2\,000--4\,000 \\
fibre launches & 4--6 $\times$ FiberPort PAF2A-5A or -7A (FC/APC) at \$674 & 2\,700--4\,000 \\
\multicolumn{3}{@{}l}{\emph{Test fibres and preparation}}\\
test cables & 8--12 $\times$ P3-375PM-FC-2 at \$366 (plus spares) & 2\,900--4\,400 \\
optional control fibres & 4 $\times$ P3-305A-FC-1 (SM300, FC/APC) at \$131 & 530 \\
polishing consumables & diamond films 30--1\,\textmu m, \SI{0.01}{\micro m} final film, FC/APC puck & $\approx$\,500 \\
\multicolumn{3}{@{}l}{\emph{Detection and logging}}\\
reference power meter & S120VC sensor + PM100USB & 1\,061 \\
arm photodiodes and DAQ & 12--18 UV-resistant photodiodes, transimpedance amplifiers, multichannel ADC & $\approx$\,1\,500--3\,000 \\
optional integrating sphere & S142C for the output detector of a key arm & 1\,215 \\
\multicolumn{3}{@{}l}{\emph{End-face cells}}\\
windows & 8--12 $\times$ WG41010-A (UV fused silica, \SI{1}{mm}, AR \SIrange{350}{700}{nm}) at \$123 & 990--1\,480 \\
FC/APC receptacles, sleeves & ADAFCPM2 (\$62) or equivalent, 4--8 & 250--500 \\
printing materials, O-rings, epoxy, inserts, PTFE membranes, desiccant & PETG, FKM/butyl O-rings, low-outgassing epoxy & $\approx$\,500--1\,000 \\
\multicolumn{3}{@{}l}{\emph{Atmosphere control and verification}}\\
gas handling & regulators, \Oxy/\water{} purifier, \SI{3}{nm} filters, flow meters or mass-flow controllers, stainless lines and fittings & $\approx$\,3\,000--6\,000 \\
trace \Oxy{} sensing & reader with ultra-trace sensor spots (e.g.\ PreSens SP-PSt9~\cite{PreSensPSt9}); optionally a galvanic trace analyser for exhausts & $\approx$\,3\,000--8\,000 \\
absorbers and indicators & RP-A and AGELESS sachets, indicators & $\approx$\,200 \\
gases & synthetic air ($\leq\SI{0.1}{ppm}$ hydrocarbons), small \Ntwo{} 6.0 cylinder, optional mixtures, rental & $\approx$\,500--1\,500/yr \\
\multicolumn{3}{@{}l}{\emph{Analysis}}\\
VOC analysis & sorbent tubes and TD-GC/MS analyses (service) & $\approx$\,500--1\,500 \\
surface analysis & SEM/EDX, AFM, Raman, ToF-SIMS/XPS facility time & $\approx$\,1\,000--4\,000 \\
\midrule
\textbf{Total if everything is bought new} & & $\approx$\,\textbf{30\,000--50\,000} \\
\end{longtable}

Using the group's existing optics, power meters, \Ntwo{} supply, polishing kit and printers roughly halves
the total. If an existing \SI{369.5}{nm} laser can spare \SIrange{20}{40}{mW} for long runs, the dedicated
source (about \$8--9k) can be dropped, at the price of competing with the experiments for laser time.
The minimum viable thesis (\cref{sec:workplan}) needs little beyond test cables, windows, printing
materials, gas lines and trace \Oxy{} verification.

\subsection{People and facilities}
Daily supervision by a group member experienced with UV fibre coupling and polishing. The mechanical
workshop for machined sealing lands and possible aluminium cells. Materials-characterisation facilities
for surface analysis, contacted in month~2. The departmental laser and gas safety officers for the risk
assessments.

\section{Risks and mitigation}\label{sec:risks}

\begin{longtable}{@{}L{4.4cm} c c L{8.0cm}@{}}
\caption{Risk register (L = likelihood, I = impact; H/M/L). To be reviewed at each monthly
meeting.}\label{tab:risks}\\
\toprule
Risk & L & I & Mitigation \\
\midrule
\endfirsthead
\toprule
Risk & L & I & Mitigation \\
\midrule
\endhead
\bottomrule
\endfoot
Testbed does not reproduce the lab degradation, or it is too slow & M & H &
Phase~0 forensics first; increase power or intensity (D1); run one arm next to the experiments
to compare environments; treat the difference as a result. \\
Alignment drift masks end-face changes & H & M &
Daily re-peak protocol with logged before/after values; stable mounts; temperature-stable enclosure
around the testbed; back-reflection and output-image monitors; sentinel arm. \\
Detector or reference drift (UV ageing of Si photodiodes, temperature) & M & M &
Low intensity on detectors, integrating spheres or diffusers, UV-resistant photodiodes, weekly cross-check with a
reference power meter, temperature logging and correction. \\
Laser instability, mode hops, shutdowns & M & L &
Ratio measurement $P_\text{out}/P_\text{ref}$; uninterruptible supply for laser and logger; automatic
restart and alert e-mail. \\
Cells leak, or printed walls are permeable & M & M &
Pressure-decay test of every cell; purge mode as fallback; epoxy sealing; metal or nylon gas lines;
continuous \Oxy{} monitoring. \\
Printed parts, absorber or sealants contaminate the end face & M & H &
\arm{N2+P} test arm; low-outgassing materials and bake-out; absorber behind a membrane; no silicone
grease; machined aluminium cells as fallback. \\
Fibre-to-fibre variability (polish, epoxy) larger than the atmosphere effect & M & H &
Single batch, one polishing procedure, acceptance criteria, randomised assignment, replicates in Campaign~II,
mixed-effects analysis. \\
Lead times or budget for parts & M & M &
Order in month~1; borrow from the group's stock; the core arms need little special equipment. \\
Access to surface-analysis facilities & M & L &
Contact facilities in month~2; keep uncleaned ends in clean, \Ntwo{}-filled containers until analysed. \\
Negative result: oxygen not the cause & M & L &
The design identifies the actual cause (\cref{tab:predictions}); this is a valuable result for the group
and for publication. \\
Student overloaded by scope & M & M &
Minimum viable thesis defined (\cref{sec:workplan}); extended arms optional; weekly priorities agreed
with the supervisor. \\
Safety incident (UV exposure, gas, chemicals) & L & H &
Risk assessments before work starts, enclosed beams, interlocks, training (\cref{sec:safety}). \\
\end{longtable}

\section{Health and safety}\label{sec:safety}

Hands-on work starts only after the departmental inductions and with signed risk assessments for the
laser work, the gas set-up, and the chemicals (COSHH). The main hazards and controls are:

\begin{description}[style=nextline,leftmargin=1.2em,itemsep=3pt]
  \item[Near-UV laser radiation] A \SI{375}{nm} diode of tens of milliwatts, or the \SI{369.5}{nm}
        systems, are Class~3B. Near-UV beams are barely visible, so stray reflections are easy to
        miss. Controls: enclose the beam paths of the testbed (the long runs are unattended), fit an
        interlock to the enclosure, wear eyewear with adequate optical density at the wavelength in
        use, use fluorescent viewing cards, register the set-up with the laser safety officer and follow
        the local rules. Unattended operation needs its own risk assessment.
  \item[Compressed and inert gases] Cylinders are secured and fitted with suitable regulators and
        relief devices. The lab \Ntwo{} supply is used where possible. The purge flows are tiny
        (about \SI{0.1}{L/min} in total) and cannot deplete the oxygen in a ventilated lab. The relevant
        scenario is the release of a whole cylinder in a small room. The risk assessment must check the
        room volume and ventilation and decide whether an oxygen-depletion monitor is needed. Exhausts are
        routed to extraction.
  \item[Forming gas] \SI{5}{\percent} \Htwo{} in \Ntwo{} is normally supplied and classified as a
        non-flammable mixture. It is still handled as a hydrogen-containing gas (no ignition sources at
        the exhaust, leak checks with a hydrogen detector). Use only if Campaign~II needs it.
  \item[Oxygen absorbers] Iron-based sachets heat up as they react. Keep them sealed until use, do not
        cut them open, and allow used sachets to finish reacting before disposal.
  \item[Chemicals] Isopropanol and acetone (flammable), epoxies and resins (skin sensitisers),
        polishing slurries. Wear gloves and eye protection and work with ventilation as stated in the
        COSHH assessments.
  \item[3D printing] Run printers in ventilated or filtered enclosures (FDM printers emit ultrafine
        particles and VOCs). Handle uncured resin with gloves.
  \item[Fibre shards] Wear eye protection when cleaving and polishing, use a sharps container and
        a dedicated work mat, and keep food and drink out of the area.
  \item[Plasma and UV/ozone cleaners] If used, they produce UV-C and ozone. Use them only with their
        interlocks and extraction.
\end{description}

\begin{safetybox}[Rule of thumb]
If an experiment is going to run unattended for weeks, design it so that every failure (power cut, gas
cylinder empty, laser fault, interlock trip) leaves it in a safe state \emph{and} is logged with a
timestamp. Then the failure costs you a data gap but does not put anyone at risk.
\end{safetybox}

\section{Expected outcomes and impact}\label{sec:outcomes}

\paragraph{For the group} Either a validated, low-cost way to stop the degradation of \SI{369.5}{nm}
fibre links (printed cells with oxygen absorbers or a clean purge), or a clear identification of the
cause and a recommendation for fibre or end-face engineering. Both reduce down time, repeated polishing,
and slow power and pointing drifts in the experiments. The same cells can protect fibres at
\SI{399}{nm} and other near-UV wavelengths used in the group.

\paragraph{For the field} Fibre delivery of near-UV light is a widespread practical problem in
trapped-ion and neutral-atom experiments and in UV photonics, but the role of the atmosphere at the end
face has, to our knowledge, not been isolated systematically. A controlled comparison of \Oxy, water,
contaminants and particles at CW near-UV intensities of order \SI{e5}{W/cm^2} is suitable for a short paper
(e.g.\ in \emph{Optics Express}, \emph{Applied Optics}, \emph{Review of Scientific Instruments} or
\emph{EPJ Techniques and Instrumentation}). Open CAD files and procedures make the result easy to adopt.

\paragraph{For the student} Training in experimental design for slow, drift-prone measurements, fibre
optics and UV laser handling, gas handling, CAD and 3D printing for functional parts, data acquisition
and statistics, and surface analysis. This mix of skills is valued in quantum-technology labs and in
industry.


\clearpage
\part*{Part II -- Guide for the student}
\addcontentsline{toc}{part}{Part II -- Guide for the student}
\section{How to use this guide}\label{sec:guide}
Part~I describes \emph{what} the project should find out and \emph{why}. This part is about
\emph{how}: what to do first, which skills to build, how to keep records that will still make sense
in month eleven, and which mistakes cost other people months. Treat it as a starting point and
improve it. If something here turns out to be wrong for our lab, write that down. Corrections like
that belong in your thesis.

\begin{keybox}[Three habits that decide the outcome]
\begin{enumerate}[leftmargin=1.2em]
  \item \textbf{Look before you touch.} Image every end face \emph{before} cleaning or polishing it. A
        cleaned-off film is evidence you can never get back.
  \item \textbf{Log every intervention.} Every re-peak, swap, cleaning, gas-cylinder change and
        power cut goes into the event log with a timestamp. Weeks-long runs are only interpretable
        if you know what happened to them.
  \item \textbf{Write as you go.} Draft the methods chapter while you build. By month~9 you will
        have forgotten why you chose a \SI{5}{mm} lens or a \SI{3}{nm} filter.
\end{enumerate}
\end{keybox}

\section{Your first four weeks}
\begin{table}[htbp]
\centering\small
\caption{Suggested onboarding plan. Adapt it with your supervisor in the first meeting.}
\label{tab:onboarding}
\begin{tabularx}{\textwidth}{@{}l Y L{4.0cm}@{}}
\toprule
Week & Activities & Output \\
\midrule
1 & Departmental, laser, gas and chemical inductions; accounts (lab notebook, repository, data storage);
    read the five starred papers in \cref{sec:reading}; watch an experienced group member polish and
    inspect a connector & reading notes (one paragraph per paper); induction certificates \\
2 & Practise cleaning, inspection and coupling with a visible laser and a spare fibre; measure
    coupling efficiency properly (\cref{sec:skills}); collect degraded fibres and their histories
    from the group & first coupling measurements; inventory of used fibres \\
3 & Phase~0 forensics on the used fibres (imaging, low-power throughput); draft risk assessments;
    first CAD model and a test print with an O-ring groove; Python logger reading one photodiode &
    image set; draft risk assessments; logger prototype \\
4 & Write the project plan (2--3 pages: aims, first experiments, parts list with prices and
    lead times); present it in the group meeting (10 minutes); place orders &
    MS1: plan and risk assessments signed off, orders placed \\
\bottomrule
\end{tabularx}
\end{table}

\section{Core skills}\label{sec:skills}

\subsection{Handling, cleaning and inspecting fibre end faces}
\begin{itemize}
  \item Never touch an end face, and cap every connector as soon as it is unplugged. Store fibres with
        caps on, in a clean box.
  \item Always \emph{inspect first}. If the face is clean, do not clean it. Cleaning tools can leave
        residues and scratch the face.
  \item For cleaning, go from dry to wet: filtered gas or a dry cleaning cassette first, then a
        lint-free wipe moistened with reagent-grade isopropanol, followed immediately by a dry wipe.
        Inspect again.
  \item Use the same microscope, magnification, illumination and focus for every image and save the
        raw files with the fibre ID, end (A/B), date and step in the file name. Use the core and cladding
        zones of IEC~61300-3-35 as the vocabulary for describing defects.
\end{itemize}

\subsection{Polishing}
Learn the group's procedure from an experienced member and then write it down as an SOP with all
parameters (films, grit sequence, times, pressure, lubricant, rinsing). Keep a polishing log for every
end. Variability between polishes is one of the largest uncontrolled factors in this project. Make
polishing boring and reproducible, and measure its outcome (inspection images, coupling efficiency,
end-face geometry if an interferometer is available).

\subsection{Coupling near-UV light into a single-mode fibre}
\begin{itemize}
  \item \textbf{Mode matching.} For a collimated Gaussian beam of $1/e^2$ radius $w_\text{in}$, a lens of
        focal length $f$ produces a waist $w_0=\lambda f/(\pi w_\text{in})$. Choose $f$ so that $w_0$
        equals the fibre mode-field radius. With a mode-field diameter of \SI{\approx2.4}{\micro m} at
        \SI{369.5}{nm}, an $f=\SI{5}{mm}$ asphere needs a \SI{\approx1.0}{mm} diameter beam
        (\cref{app:estimates}).
  \item \textbf{Alignment.} Start with a visible alignment beam if the optics allow it. Make the
        output of the fibre visible, then walk the beam with two mirrors (translation and angle) before
        touching focus. Optimise the focus last and re-walk.
  \item \textbf{PM fibre.} Align the input polarisation to the slow axis: rotate the input half-wave
        plate while measuring the output polarisation extinction ratio, and confirm by gently warming or
        stressing the fibre (the output polarisation should not change).
  \item \textbf{Measure coupling efficiency correctly.} Measure the power directly before the fibre
        (after all optics, including the cell window) and directly after the output, with the same
        calibrated sensor. Correct for the Fresnel loss at each face if you report an ``internal''
        efficiency, and state clearly which definition you use.
\end{itemize}

\subsection{Measuring UV power reliably}
\begin{itemize}
  \item Silicon photodiodes slowly lose responsivity under UV exposure, and photodiodes and amplifiers
        drift with temperature. Keep the intensity on detectors low, prefer integrating spheres or
        diffusers, and consider photodiodes specified for UV resistance for the long runs. (SiC photodiodes,
        often recommended for UV, barely respond at \SI{369.5}{nm}.)
  \item Check the zero (laser off) and the linearity (calibrated attenuator) at the start of every run.
  \item Avoid interference fringes: use wedged optics and tilt surfaces slightly, otherwise
        etalons turn tiny temperature drifts into apparent power changes.
  \item Cross-check all detectors weekly against one reference power meter that is otherwise kept in a
        drawer (so it does not age).
\end{itemize}

\subsection{Gas handling and leak testing}
\begin{itemize}
  \item Get trained on regulators and cylinder handling. Use compression fittings correctly (do not
        over-tighten; replace ferrules) and leak-check every new joint with a leak-detection liquid or a
        helium or hydrogen sniffer.
  \item For ppm-level oxygen, the tubing matters. Polyurethane and fluoropolymer tubes are permeable
        to \Oxy; stainless steel, copper or short nylon lines are better (\cref{app:estimates}).
  \item Leak-test every cell before use: pressurise to about \SI{+100}{mbar} with \Ntwo, log pressure
        \emph{and} temperature for 12 hours, and compute the leak rate $Q=V\,\mathrm{d}p/\mathrm{d}t$ after
        correcting for temperature (\cref{app:sops}).
\end{itemize}

\subsection{CAD and 3D printing of gas-tight parts}
\begin{itemize}
  \item FDM prints are porous unless printed with \SI{100}{\percent} infill, several perimeters, slight
        over-extrusion and good layer adhesion. Even then, seal the inside with a thin low-outgassing epoxy
        for the static cells. Resin (SLA) prints are denser but release residual monomer unless fully
        post-cured. Bake parts (at a temperature safely below the glass transition of the polymer) under
        \Ntwo{} purge before use.
  \item Design O-ring grooves to standard proportions (squeeze and gland fill as in a sealing handbook)
        and print the sealing faces flat and smooth, or face them with a thin epoxy layer. Use heat-set
        metal inserts for screws.
  \item Choose the material with \cref{tab:materials} in mind. PETG is a good first choice for FDM. ABS
        (styrene emission, high water permeability) and standard SLA resins (high outgassing) are poor
        choices near UV-illuminated end faces. Use walls of at least \SI{1.1}{mm} and randomised
        seams~\cite{Gordeev2018}.
  \item Print a test coupon with every new material or setting, and keep coupons for the \arm{N2+P}
        outgassing test.
\end{itemize}

\subsection{Data acquisition and analysis}
\begin{itemize}
  \item Write the logger so that it never loses data. Write to a new file every day, flush after each
        write, use UTC timestamps, and save a metadata file with fibre IDs, arm assignment, calibration
        factors and the software version (git commit). Make it restart automatically after a power cut.
  \item Keep raw data read-only. All corrections (detector temperature, flagged intervals) are applied
        by analysis scripts, never by editing data files.
  \item Make every figure in the thesis reproducible from raw data with one script. A notebook is fine
        for exploration, but the final analysis should be a script in the repository.
  \item Learn enough statistics to defend an upper limit: regression with autocorrelated residuals,
        block bootstrap, and mixed-effects models (\cref{sec:analysis}).
\end{itemize}

\section{Working practices}
\begin{description}[style=nextline,leftmargin=1.2em,itemsep=3pt]
  \item[Lab notebook] Use the group's electronic notebook. Create one entry per day with the plan,
        what was done, observations, and links to data and images. Note surprises immediately, even if
        they seem unimportant.
  \item[Repository] Keep code, CAD files, SOPs, and the thesis source in version control. Commit small
        and often with meaningful messages. Store large raw data on the group's storage, not in git, and
        put the paths in the metadata.
  \item[Meetings] Meet the daily supervisor weekly for 30 minutes with a short written agenda (done,
        problems, plan, decisions needed). Meet the PI monthly, together with a one-page progress
        report that also reviews the risk register (\cref{tab:risks}). Present in the group meeting
        about every two months.
  \item[Asking for help] If you are stuck for more than a day on something that someone in the group
        has done before, ask. Experienced colleagues can often solve fibre-coupling and vacuum or gas
        problems in minutes.
\end{description}

\section{Common pitfalls}
\begin{pitfallbox}[Lessons from long-term optical measurements]
\begin{itemize}[leftmargin=1em]
  \item \textbf{Confusing alignment drift with degradation.} Always record throughput before
        \emph{and} after re-peaking.
  \item \textbf{Cleaning away the evidence.} Image and, if possible, analyse before cleaning.
  \item \textbf{Comparing runs done at different times.} Run conditions in parallel; if you must
        compare runs, keep a reference arm common to both.
  \item \textbf{Trusting a gas cell you have not measured.} A cell with an indicator that has turned
        colour, or an \Oxy{} reading that drifts up, means the run is compromised. Log \Oxy{}, RH and
        flow continuously.
  \item \textbf{Introducing contamination with the fix.} Silicone grease, fresh resin prints, adhesive tapes
        and cable jackets near the end face can all release organics. Keep them out of the cell or
        test them (\arm{N2+P}).
  \item \textbf{Changing two things at once.} New fibre batch, new polish and new atmosphere in the same
        run cannot be untangled.
  \item \textbf{Leaving the thesis until the end.} Writing always takes longer than planned, and writing
        early also shows where data are missing while you can still take them.
\end{itemize}
\end{pitfallbox}

\section{Writing the thesis}
A suggested structure (adapt it to the departmental regulations):
\begin{enumerate}
  \item Introduction: UV light delivery for trapped ions, the problem, aims.
  \item Background: optical fibres and coupling; defects and solarisation in silica; surface
        photochemistry and laser-induced contamination; gas transport and controlled atmospheres.
  \item Methods: testbed, cells and their qualification, procedures, data analysis.
  \item Results I: forensics of used fibres and the baseline degradation law.
  \item Results II: the atmosphere campaigns and the cleaning sequence; surface analysis.
  \item Results III: demonstration on a \SI{369.5}{nm} link.
  \item Discussion: mechanism, comparison with the literature, design guidelines.
  \item Conclusions and outlook.
\end{enumerate}
Appendices: CAD drawings and bill of materials, SOPs, logger documentation, data-availability statement.
Plan for at least eight weeks of focused writing, and give your supervisor complete chapter drafts
early enough for a round of feedback.

\section{Reading list}\label{sec:reading}
Read the starred items in the first week, the rest as the project needs them. For every paper, note in
your reading log the wavelength, intensity, atmosphere, fibre or surface type, time scale and the
numbers you might want to compare with later. A one-page table of these numbers will be the backbone
of the discussion chapter.

\subsection*{Essential (week 1)}
\begin{enumerate}[label=$\bigstar$,leftmargin=1.6em]
  \item Colombe et al.~\cite{Colombe2014}: the reference paper on UV solarisation of single-mode
        fibres for ion trapping, and on hydrogen loading and curing.
  \item Marciniak et al.~\cite{Marciniak2017}: UV patch cords in an ion-trap lab, facet deposits and
        the nitrogen purge. Read the methods carefully: they are close to what you will build.
  \item Gonschior et al.~\cite{Gonschior2013JNCS} and the summary chapters of Gonschior's
        thesis~\cite{Gonschior2013PhD}: the most detailed study of CW blue-light degradation of
        small-core single-mode fibres, including end caps.
  \item Mondin et al.~\cite{Mondin2017} and Hippler et al.~\cite{Hippler2016}: what the space-laser
        community learned about laser-induced contamination and oxygen.
  \item Parthasarathy et al.~\cite{Parthasarathy2024}: CW contamination at an intensity similar to
        ours, as a function of oxygen pressure.
\end{enumerate}

\subsection*{Background for the thesis (weeks 2--8)}
\begin{itemize}
  \item Defects in silica: Skuja~\cite{Skuja1998}, Skuja et al.~\cite{Skuja2005}, Griscom~\cite{Griscom1991}
        and the fibre-oriented review by Girard et al.~\cite{Girard2019}. Oxygen species in silica:
        Kajihara et al.~\cite{Kajihara2008a}.
  \item Surfaces: Radzig~\cite{Radzig2000} on activated silica surfaces.
  \item Fibre alternatives: Gebert et al.~\cite{Gebert2014} and Dorer et al.~\cite{Dorer2024} (hollow core),
        Yamamoto et al.~\cite{Yamamoto2009} (LMA PCF in the deep UV).
  \item Oxygen in UV laser systems: Wilson et al.~\cite{Wilson2011} and Hannig et al.~\cite{Hannig2018}
        (UV conversion crystals in ion-trap labs), Schoedl et al.~\cite{Schoedl2005} (laser-diode facets in
        nitrogen versus oxygen), and the Fujifilm patent on siloxane deposits at fibre end
        faces~\cite{Fujifilmpatent2006}.
  \item Optics: Kogelnik and Li~\cite{Kogelnik1966} on Gaussian beams (mode matching), Lax~\cite{Lax1977}
        on laser heating.
  \item Cleaning: Vig~\cite{Vig1985} on UV/ozone cleaning.
\end{itemize}

\subsection*{Practical references}
IEC 61300-3-35 for end-face inspection~\cite{IEC61300335}; Thorlabs' polishing instructions and fibre
damage-threshold tutorial~\cite{ThorlabsPolish,ThorlabsDamage}; the NIST UV patch-cord
recipe~\cite{NISTrecipe}; photodiode data sheets with notes on UV exposure~\cite{HamamatsuS1336,HamamatsuS12698};
and the references on 3D-printed vacuum and gas parts in \cref{sec:sota-atmosphere}.


\clearpage
\appendix
\part*{Appendices}
\addcontentsline{toc}{part}{Appendices}
\section{Placeholder 14-estimates}\label{sec:ph-14-estimates}
\label{app:estimates}

\section{Draft standard operating procedures}\label{app:sops}
These drafts are a starting point. The student turns them into the group's SOPs during WP0--WP3, with
the actual equipment, part numbers and parameters.

\subsection*{SOP-1: End-face inspection and cleaning}
\begin{enumerate}
  \item Wear gloves; work on a clean mat; have caps ready.
  \item Inspect the end face with the inspection microscope at the standard settings. Save the image as
        \texttt{<fibreID>\_<end>\_<YYYYMMDD>\_<step>.png}.
  \item If clean (no defects in the core zone, nothing loose in the cladding zone): stop and cap.
  \item Dry clean (filtered gas, then a dry cleaning cassette or click cleaner); inspect and image.
  \item If still contaminated: wet--dry clean with reagent-grade isopropanol on a lint-free wipe, then
        immediately a dry wipe; inspect and image. Repeat at most twice.
  \item Record every step in the fibre's log. Never clean a fibre that is scheduled for surface analysis.
\end{enumerate}

\subsection*{SOP-2: Polishing and acceptance of a test end}
\begin{enumerate}
  \item Polish with the group's procedure (films, sequence, times, pressure, lubricant written down) and
        log all parameters.
  \item Clean and inspect (SOP-1). Measure end-face geometry if an interferometer is available.
  \item Measure the coupling efficiency at low power (\SI{\le 0.3}{mW}) with the standard launch.
  \item Accept if the core zone is defect-free and the efficiency is within the agreed band around the
        batch median; otherwise re-polish. Cap and store the fibre in a clean container until use.
\end{enumerate}

\subsection*{SOP-3: Cell qualification (leak and oxygen hold test)}
\begin{enumerate}
  \item Clean all parts (isopropanol wipe, no silicone grease), bake printed parts under \Ntwo{} purge,
        assemble with a dummy connector.
  \item Pressure decay: pressurise to \SI{+100}{mbar} with \Ntwo, close the valves, log pressure and
        temperature for \SI{12}{h}. Correct for temperature ($\Delta p/p=\Delta T/T$) and compute
        $Q=V\,\mathrm{d}p/\mathrm{d}t$. Accept if $Q<\SI{1e-5}{mbar.L/s}$.
  \item Oxygen hold (static cells): purge with \Ntwo{} to $\lesssim\SI{1}{\percent}$ \Oxy, insert the
        absorber, seal, and log the \Oxy{} sensor spot reading for at least \SI{7}{days}. Accept if \Oxy{}
        stays below \SI{100}{ppm} after the initial absorption period.
  \item Record the cell ID, material, print settings, sealing and test results in the cell log.
\end{enumerate}

\subsection*{SOP-4: Start of a run}
\begin{enumerate}
  \item Fibres accepted (SOP-2), IDs and ends assigned to arms by randomisation, assignment saved in the
        run metadata.
  \item Cells qualified (SOP-3); purge flows set and logged; \Oxy, RH and VOC readings stable.
  \item Detectors zeroed; reference power meter cross-check recorded.
  \item Couple each arm at low power, then raise to the test power and set equal powers; record the
        initial values.
  \item Start the logger; write the run plan (duration, re-peak schedule, stop criteria) into the notebook.
\end{enumerate}

\subsection*{SOP-5: End of a run}
\begin{enumerate}
  \item Record the final throughput before and after re-peaking at the test power, then at low power.
  \item Switch off the laser. Inspect both ends without touching (SOP-1, imaging only).
  \item Either run the cleaning sequence of \cref{sec:ladder} (output end first) with a low-power
        measurement after each step, or cap the fibre and store it in an \Ntwo-filled, labelled container
        for surface analysis.
  \item Archive the data, images and notebook entries of the run and write a one-page run summary.
\end{enumerate}

\section{Data-logging specification}\label{app:logging}
\begin{table}[h]
\centering\small
\caption{Minimum channel list for the logger. Raw sampling at \SI{1}{Hz}; stored as 1-minute means
with standard deviations; one file per day plus a metadata file per run.}
\label{tab:channels}
\begin{tabularx}{\textwidth}{@{}l Y l@{}}
\toprule
Channel & Description & Units \\
\midrule
\texttt{t\_utc} & timestamp (ISO 8601, UTC) & -- \\
\texttt{arm<k>\_ref}, \texttt{arm<k>\_out}, \texttt{arm<k>\_back} & photodiode voltages of arm $k$ & V \\
\texttt{arm<k>\_O2}, \texttt{arm<k>\_RH}, \texttt{arm<k>\_VOC} & atmosphere of arm $k$ (cell or outlet) & ppm, \%, index \\
\texttt{arm<k>\_flow}, \texttt{arm<k>\_T} & purge flow, temperature at the cell & sccm, \si{\celsius} \\
\texttt{laser\_I}, \texttt{laser\_T}, \texttt{laser\_P} & diode current, temperature, monitor & mA, \si{\celsius}, V \\
\texttt{lab\_T}, \texttt{lab\_RH}, \texttt{lab\_p} & lab environment & \si{\celsius}, \%, hPa \\
\texttt{event} & free-text event log entries (re-peak, swap, cleaning, outage) & -- \\
\bottomrule
\end{tabularx}
\end{table}
The run metadata (JSON or YAML) contains: run ID, start and end time, fibre IDs and ends per arm,
atmosphere per arm, gas cylinder IDs and certificates, cell IDs, test power, detector calibration
factors, software version (git commit) and the analysis plan version.


\clearpage
\printbibliography[heading=bibintoc,title={References}]

\end{document}
